\documentclass[11pt,a4paper]{article}

\usepackage[utf8]{inputenc}
\usepackage[T1]{fontenc}
\usepackage[french]{babel}

\usepackage{geometry}
\geometry{margin=2.5cm}

\usepackage{hyperref}
\usepackage{enumitem}
\usepackage{graphicx}
\usepackage{amsmath}
\usepackage{tcolorbox}

\title{Cahier des Charges\\
Agent Personnel de Recherche de Stage}

\author{Projet personnel}

\date{\today}

\begin{document}

\maketitle

\tableofcontents

\newpage

\section{Introduction}

Ce document constitue le cahier des charges initial d’un système logiciel visant à automatiser et assister la recherche de stage à l’aide de techniques d’intelligence artificielle et d’automatisation.

L'objectif est de concevoir un \textbf{agent personnel de recherche de stage} capable de surveiller en continu différentes sources d’offres, d’analyser leur pertinence par rapport au profil de l’utilisateur, de préparer les candidatures, puis d’assister leur soumission et leur suivi.

Le système doit agir de manière proactive tout en conservant une \textbf{validation humaine obligatoire} avant toute action externe telle que l’envoi d’une candidature ou d’une relance.

\section{Objectifs du système}

L’objectif principal du système est de reproduire le comportement d’un utilisateur dans sa recherche de stage tout en automatisant les tâches répétitives.

Plus précisément, le système doit :

\begin{itemize}
\item surveiller régulièrement les plateformes proposant des stages ;
\item détecter les nouvelles offres pertinentes ;
\item analyser automatiquement la compatibilité entre une offre et le profil de l’utilisateur ;
\item préparer les éléments de candidature (CV, lettre, formulaires) ;
\item assister la soumission de candidature avec validation humaine ;
\item assurer le suivi des candidatures envoyées ;
\item proposer et préparer des relances en cas d’absence de réponse.
\end{itemize}

\section{Principe général}

Le système doit fonctionner comme un \textbf{assistant personnel autonome} reproduisant le comportement d’un chercheur de stage.

Le cycle de fonctionnement global est le suivant :

\begin{enumerate}
\item surveillance périodique des sources d'offres ;
\item détection des nouvelles offres ;
\item filtrage selon les préférences configurées ;
\item analyse de correspondance avec le profil utilisateur ;
\item préparation des candidatures pertinentes ;
\item validation humaine avant envoi ;
\item suivi des candidatures ;
\item relances éventuelles.
\end{enumerate}

\section{Sources d’offres}

Le système doit être capable d’explorer différentes sources d’offres de stage, notamment :

\begin{itemize}
\item LinkedIn
\item Indeed
\item Welcome to the Jungle
\item plateformes d’emploi spécialisées
\item pages carrières des entreprises
\item career centers d’écoles ou d’universités
\end{itemize}

Cette liste pourra être enrichie au cours du développement.

\section{Cycle de fonctionnement}

Le système fonctionne selon un cycle déclenché périodiquement (par exemple toutes les trois heures).

Lors de chaque cycle, les étapes suivantes sont exécutées.

\subsection{Collecte des offres}

Le système visite les sources configurées et collecte les nouvelles offres disponibles.

\subsection{Détection des nouvelles offres}

Les offres collectées sont comparées avec celles déjà enregistrées dans la base afin de détecter :

\begin{itemize}
\item les nouvelles offres
\item les offres déjà connues
\item les offres mises à jour
\end{itemize}

\subsection{Filtrage initial}

Un filtrage est appliqué selon les critères configurés par l’utilisateur :

\begin{itemize}
\item type de contrat
\item durée
\item localisation
\item mots-clés
\item entreprises cibles
\item entreprises exclues
\end{itemize}

\subsection{Analyse de correspondance}

Pour chaque offre retenue, le système analyse la correspondance entre l’offre et le profil utilisateur.

Cette analyse produit :

\begin{itemize}
\item un score global de correspondance
\item une justification détaillée
\item une estimation de priorité
\end{itemize}

\subsection{Décision}

Selon le score obtenu, le système peut :

\begin{itemize}
\item ignorer l’offre
\item la notifier à l’utilisateur
\item préparer une candidature
\end{itemize}

\section{Préparation des candidatures}

Pour les offres jugées pertinentes, le système doit être capable de :

\begin{itemize}
\item sélectionner la version la plus appropriée du CV
\item adapter certains éléments du CV si nécessaire
\item générer une lettre de motivation personnalisée
\item préparer les informations nécessaires aux formulaires
\end{itemize}

Les candidatures préparées sont placées dans un état \textbf{prêtes à soumettre}.

\section{Validation humaine}

Aucune candidature ne doit être envoyée sans validation explicite de l’utilisateur.

Le système doit permettre à l’utilisateur de :

\begin{itemize}
\item valider la candidature
\item modifier les documents générés
\item rejeter la candidature
\item reporter la décision
\end{itemize}

\section{Suivi des candidatures}

Après soumission, le système doit enregistrer les informations suivantes :

\begin{itemize}
\item entreprise
\item poste
\item date de candidature
\item canal de candidature
\item statut de la candidature
\end{itemize}

Le système doit maintenir un historique complet des interactions.

\section{Gestion des réponses}

Le système peut être connecté à une messagerie (par exemple Gmail) afin de :

\begin{itemize}
\item détecter les réponses des recruteurs
\item associer les réponses aux candidatures correspondantes
\item notifier l’utilisateur
\item proposer des réponses adaptées
\end{itemize}

\section{Gestion des relances}

En l’absence de réponse dans un délai configuré, le système peut préparer une relance.

Les relances doivent :

\begin{itemize}
\item être rédigées automatiquement
\item être adaptées au contexte de la candidature
\item nécessiter une validation humaine avant envoi
\end{itemize}

\section{Contraintes du système}

Le système doit respecter les contraintes suivantes :

\begin{itemize}
\item ne jamais envoyer de candidature sans validation humaine ;
\item ne jamais modifier les informations du CV de manière mensongère ;
\item éviter les candidatures en double ;
\item conserver un historique complet des actions ;
\item permettre l’audit et la traçabilité des décisions.
\end{itemize}

\section{Évolution du document}

Ce cahier des charges constitue une première version du projet.

Il sera enrichi progressivement avec :

\begin{itemize}
\item l’architecture technique
\item la modélisation des données
\item la description des modules logiciels
\item les choix technologiques
\item les scénarios d’utilisation détaillés
\end{itemize}

\section{Modélisation fonctionnelle du système}

\subsection{Acteurs du système}

Le système met en jeu plusieurs acteurs.

\begin{itemize}
    \item \textbf{Utilisateur} : configure le système, valide les candidatures, corrige les documents générés, supervise le suivi.
    \item \textbf{Moteur de veille} : collecte périodiquement les offres depuis les sources configurées.
    \item \textbf{Moteur de décision} : filtre, analyse et score les offres selon le profil utilisateur.
    \item \textbf{Moteur de préparation} : prépare les documents de candidature et les informations nécessaires à la soumission.
    \item \textbf{Moteur de suivi} : surveille les réponses, met à jour les statuts et prépare les relances.
    \item \textbf{Services externes} : plateformes d’offres, messagerie, services de notification.
\end{itemize}

\subsection{Objets métier}

Les principaux objets manipulés par le système sont les suivants :

\begin{itemize}
    \item \textbf{Offre} : annonce de stage détectée sur une source.
    \item \textbf{Source} : plateforme ou site à partir duquel les offres sont collectées.
    \item \textbf{Entreprise} : organisation proposant une offre.
    \item \textbf{Profil utilisateur} : ensemble des préférences, compétences et contraintes de l’utilisateur.
    \item \textbf{Version de CV} : version spécifique du CV disponible pour adaptation ou soumission.
    \item \textbf{Brouillon de candidature} : ensemble préparatoire avant validation.
    \item \textbf{Candidature} : candidature effectivement soumise.
    \item \textbf{Relance} : message de suivi préparé après absence de réponse.
    \item \textbf{Réponse} : message reçu d’un recruteur ou d’une entreprise.
    \item \textbf{Journal de décision} : enregistrement des décisions prises par le système.
\end{itemize}

\subsection{États d’une offre}

Une offre peut évoluer selon les états suivants :

\begin{center}
DETECTEE $\rightarrow$ FILTREE $\rightarrow$ ANALYSEE $\rightarrow$ QUALIFIEE
$\rightarrow$ A\_PREPARER $\rightarrow$ PREPAREE $\rightarrow$ A\_VALIDER
$\rightarrow$ VALIDEE $\rightarrow$ SOUMISE $\rightarrow$ EN\_SUIVI
$\rightarrow$ CLOTUREE
\end{center}

Des branches alternatives sont possibles :

\begin{itemize}
    \item ANALYSEE $\rightarrow$ REJETEE
    \item PREPAREE $\rightarrow$ ABANDONNEE
    \item A\_VALIDER $\rightarrow$ MODIFICATION\_DEMANDEE
    \item EN\_SUIVI $\rightarrow$ REFUS
    \item EN\_SUIVI $\rightarrow$ ENTRETIEN
    \item EN\_SUIVI $\rightarrow$ ACCEPTATION
\end{itemize}

\subsection{États d’une relance}

Le processus de relance suit les états suivants :

\begin{center}
AUCUNE\_RELANCE $\rightarrow$ RELANCE\_A\_PREPARER
$\rightarrow$ RELANCE\_A\_VALIDER
$\rightarrow$ RELANCE\_ENVOYEE
\end{center}

En cas de retour recruteur :

\begin{center}
RELANCE\_ENVOYEE $\rightarrow$ REPONSE\_RECUE
\end{center}

\subsection{Cas d’usage principaux}

Les cas d’usage principaux du système sont les suivants :

\begin{enumerate}
    \item détecter une nouvelle offre depuis une source configurée ;
    \item filtrer l’offre selon les préférences utilisateur ;
    \item analyser la pertinence de l’offre par rapport au profil ;
    \item préparer une candidature pour une offre qualifiée ;
    \item soumettre la candidature à validation humaine ;
    \item enregistrer la candidature après envoi ;
    \item surveiller les réponses reçues via les canaux connectés ;
    \item préparer une relance en l’absence de réponse ;
    \item clôturer une candidature selon son issue.
\end{enumerate}

\section{Modèle de données conceptuel}

Cette section décrit les principales entités manipulées par le système ainsi que leurs relations conceptuelles.

L’objectif est de structurer les informations nécessaires au bon fonctionnement du produit avant la définition détaillée du schéma physique de base de données.

\subsection{Principes généraux}

Le modèle de données doit permettre :

\begin{itemize}
    \item le stockage des offres détectées ;
    \item la conservation des sources d’origine ;
    \item l’enregistrement du profil utilisateur et des préférences ;
    \item la gestion des différentes versions de CV et des lettres générées ;
    \item le suivi complet des candidatures ;
    \item la préparation et le suivi des relances ;
    \item la traçabilité des décisions automatiques et humaines ;
    \item l’historisation des exécutions du système.
\end{itemize}

\subsection{Entités principales}

\subsubsection{Source}

L’entité \textbf{Source} représente une plateforme ou un canal à partir duquel des offres sont collectées.

Exemples :
\begin{itemize}
    \item LinkedIn ;
    \item Indeed ;
    \item page carrière d’une entreprise ;
    \item career center ;
    \item flux email ou newsletter.
\end{itemize}

Attributs conceptuels possibles :
\begin{itemize}
    \item identifiant de la source ;
    \item nom ;
    \item type de source ;
    \item URL de base ;
    \item statut d’activation ;
    \item fréquence de surveillance.
\end{itemize}

\subsubsection{Entreprise}

L’entité \textbf{Entreprise} représente l’organisation qui publie ou porte une offre.

Attributs conceptuels possibles :
\begin{itemize}
    \item identifiant de l’entreprise ;
    \item nom ;
    \item secteur ;
    \item site web ;
    \item localisation principale ;
    \item niveau de priorité ou d’intérêt ;
    \item statut cible / exclue / neutre.
\end{itemize}

\subsubsection{Offre brute}

L’entité \textbf{OffreBrute} représente l’offre telle qu’elle a été collectée depuis une source, avant normalisation complète.

Elle permet de conserver une trace fidèle de la donnée d’origine.

Attributs conceptuels possibles :
\begin{itemize}
    \item identifiant ;
    \item source associée ;
    \item identifiant externe sur la plateforme ;
    \item URL de l’offre ;
    \item titre brut ;
    \item contenu brut ;
    \item date de collecte ;
    \item date de publication détectée ;
    \item statut de parsing.
\end{itemize}

\subsubsection{Offre normalisée}

L’entité \textbf{Offre} représente une offre après nettoyage, extraction des champs utiles et dédoublonnage.

Attributs conceptuels possibles :
\begin{itemize}
    \item identifiant ;
    \item entreprise associée ;
    \item source principale ;
    \item titre normalisé ;
    \item description normalisée ;
    \item type de contrat ;
    \item durée ;
    \item localisation ;
    \item mode de travail (présentiel, hybride, télétravail) ;
    \item date de publication ;
    \item date limite éventuelle ;
    \item niveau d’étude requis ;
    \item statut courant ;
    \item score global ;
    \item score d’action ;
    \item justification de scoring.
\end{itemize}

\subsubsection{Profil utilisateur}

L’entité \textbf{ProfilUtilisateur} représente le profil métier utilisé pour l’analyse des offres.

Attributs conceptuels possibles :
\begin{itemize}
    \item identifiant ;
    \item nom du profil ;
    \item formation ;
    \item compétences ;
    \item expériences ;
    \item projets ;
    \item préférences de localisation ;
    \item types de postes recherchés ;
    \item mots-clés prioritaires ;
    \item contraintes ;
    \item date de mise à jour.
\end{itemize}

\subsubsection{Préférences de recherche}

L’entité \textbf{PreferenceRecherche} regroupe les règles de filtrage et de priorisation utilisées par le système.

Attributs conceptuels possibles :
\begin{itemize}
    \item identifiant ;
    \item fréquence de surveillance ;
    \item liste de localisations autorisées ;
    \item liste d’entreprises cibles ;
    \item liste d’entreprises exclues ;
    \item durée minimale et maximale ;
    \item date de début souhaitée ;
    \item seuil minimal de score ;
    \item politique de relance.
\end{itemize}

\subsubsection{Version de CV}

L’entité \textbf{CVVersion} représente une version de CV disponible pour la préparation des candidatures.

Attributs conceptuels possibles :
\begin{itemize}
    \item identifiant ;
    \item nom de version ;
    \item fichier associé ;
    \item orientation métier ;
    \item date de création ;
    \item version active ou non ;
    \item commentaires.
\end{itemize}

\subsubsection{Lettre de motivation}

L’entité \textbf{LettreMotivation} représente une lettre générée ou retravaillée pour une offre donnée.

Attributs conceptuels possibles :
\begin{itemize}
    \item identifiant ;
    \item offre associée ;
    \item contenu ;
    \item statut ;
    \item date de génération ;
    \item date de validation ;
    \item version finale.
\end{itemize}

\subsubsection{Brouillon de candidature}

L’entité \textbf{BrouillonCandidature} représente un dossier préparatoire avant validation humaine.

Attributs conceptuels possibles :
\begin{itemize}
    \item identifiant ;
    \item offre associée ;
    \item CV sélectionné ;
    \item lettre associée ;
    \item email ou message préparé ;
    \item réponses préparées aux questions éventuelles ;
    \item statut du brouillon ;
    \item date de création ;
    \item date de dernière modification.
\end{itemize}

\subsubsection{Candidature}

L’entité \textbf{Candidature} représente une candidature effectivement soumise.

Attributs conceptuels possibles :
\begin{itemize}
    \item identifiant ;
    \item offre associée ;
    \item brouillon d’origine ;
    \item canal de soumission ;
    \item date d’envoi ;
    \item statut courant ;
    \item date du dernier changement de statut ;
    \item commentaires utilisateur.
\end{itemize}

\subsubsection{Relance}

L’entité \textbf{Relance} représente une relance préparée ou envoyée à la suite d’une candidature.

Attributs conceptuels possibles :
\begin{itemize}
    \item identifiant ;
    \item candidature associée ;
    \item type de relance ;
    \item contenu ;
    \item date prévue ;
    \item date d’envoi ;
    \item statut.
\end{itemize}

\subsubsection{Réponse}

L’entité \textbf{Reponse} représente une réponse reçue d’un recruteur ou d’une entreprise.

Attributs conceptuels possibles :
\begin{itemize}
    \item identifiant ;
    \item candidature associée ;
    \item canal de réception ;
    \item expéditeur ;
    \item sujet ;
    \item contenu ;
    \item date de réception ;
    \item classification automatique.
\end{itemize}

\subsubsection{Exécution d’agent}

L’entité \textbf{AgentRun} représente une exécution planifiée ou manuelle du système.

Attributs conceptuels possibles :
\begin{itemize}
    \item identifiant ;
    \item date de démarrage ;
    \item date de fin ;
    \item type d’exécution ;
    \item nombre d’offres détectées ;
    \item nombre d’offres retenues ;
    \item statut final ;
    \item journal synthétique.
\end{itemize}

\subsubsection{Journal de décision}

L’entité \textbf{DecisionLog} représente une décision prise par le système ou par l’utilisateur.

Attributs conceptuels possibles :
\begin{itemize}
    \item identifiant ;
    \item type d’objet concerné ;
    \item identifiant de l’objet concerné ;
    \item auteur de la décision (système ou utilisateur) ;
    \item décision prise ;
    \item justification ;
    \item date ;
    \item métadonnées associées.
\end{itemize}

\subsection{Relations principales entre les entités}

Les relations conceptuelles principales sont les suivantes :

\begin{itemize}
    \item une \textbf{Source} peut produire plusieurs \textbf{OffresBrutes} ;
    \item une \textbf{OffreBrute} peut être transformée en une \textbf{Offre} ;
    \item une \textbf{Entreprise} peut être associée à plusieurs \textbf{Offres} ;
    \item une \textbf{Offre} peut donner lieu à zéro ou un \textbf{BrouillonCandidature} ;
    \item un \textbf{BrouillonCandidature} utilise une \textbf{CVVersion} et éventuellement une \textbf{LettreMotivation} ;
    \item une \textbf{Offre} peut donner lieu à zéro ou une \textbf{Candidature} ;
    \item une \textbf{Candidature} peut avoir plusieurs \textbf{Relances} ;
    \item une \textbf{Candidature} peut recevoir plusieurs \textbf{Reponses} ;
    \item une \textbf{DecisionLog} peut être rattachée à une \textbf{Offre}, un \textbf{BrouillonCandidature}, une \textbf{Candidature} ou une \textbf{Relance} ;
    \item un \textbf{AgentRun} peut produire plusieurs décisions et plusieurs objets métiers.
\end{itemize}

\subsection{Contraintes conceptuelles}

Le modèle de données doit respecter les contraintes suivantes :

\begin{itemize}
    \item une même offre ne doit pas être dupliquée inutilement après normalisation ;
    \item une candidature ne peut exister que si une offre correspondante existe ;
    \item une relance ne peut être créée que pour une candidature déjà envoyée ;
    \item toute décision importante doit être historisée ;
    \item toute soumission externe doit être traçable ;
    \item les versions de CV et de lettres doivent être conservées.
\end{itemize}

\subsection{Remarque sur la suite}

Le présent modèle conceptuel constitue la base de la future modélisation logique puis physique de la base de données.

Il sera complété dans la suite du projet par :

\begin{itemize}
    \item un schéma relationnel détaillé ;
    \item les clés primaires et étrangères ;
    \item les cardinalités précises ;
    \item les contraintes d’intégrité ;
    \item les choix d’indexation.
\end{itemize}

\section{Schéma relationnel logique}

Cette section propose une première traduction du modèle conceptuel en schéma relationnel logique.

L’objectif est de préparer la future implémentation de la base de données relationnelle du système, en identifiant les tables principales, leurs attributs structurants et les relations entre elles.

\subsection{Principes de modélisation}

Le schéma relationnel doit respecter les principes suivants :

\begin{itemize}
    \item séparation entre données brutes collectées et données normalisées ;
    \item conservation de l’historique des actions et décisions ;
    \item traçabilité complète des candidatures et relances ;
    \item extensibilité pour les futures versions du système ;
    \item compatibilité avec un système de workflow piloté par états.
\end{itemize}

\subsection{Tables principales}

\subsubsection{Table \texttt{sources}}

Cette table recense les différentes sources surveillées par le système.

\begin{center}
\begin{tabular}{|l|l|l|}
\hline
\textbf{Champ} & \textbf{Type logique} & \textbf{Description} \\
\hline
id & PK & identifiant interne \\
name & texte & nom de la source \\
source_type & texte & type de source \\
base_url & texte & URL principale \\
is_active & booléen & source active ou non \\
check_frequency_hours & entier & fréquence de surveillance \\
created_at & datetime & date de création \\
updated_at & datetime & date de mise à jour \\
\hline
\end{tabular}
\end{center}

\subsubsection{Table \texttt{companies}}

Cette table stocke les entreprises liées aux offres.

\begin{center}
\begin{tabular}{|l|l|l|}
\hline
\textbf{Champ} & \textbf{Type logique} & \textbf{Description} \\
\hline
id & PK & identifiant interne \\
name & texte & nom de l’entreprise \\
sector & texte & secteur d’activité \\
website_url & texte & site web \\
main_location & texte & localisation principale \\
priority_level & entier & niveau d’intérêt \\
company_status & texte & cible, neutre, exclue \\
created_at & datetime & date de création \\
updated_at & datetime & date de mise à jour \\
\hline
\end{tabular}
\end{center}

\subsubsection{Table \texttt{offers\_raw}}

Cette table conserve les offres telles qu’elles ont été collectées.

\begin{center}
\begin{tabular}{|l|l|l|}
\hline
\textbf{Champ} & \textbf{Type logique} & \textbf{Description} \\
\hline
id & PK & identifiant interne \\
source_id & FK & référence vers \texttt{sources} \\
external_offer_id & texte & identifiant externe éventuel \\
offer_url & texte & URL de l’offre \\
raw_title & texte & titre brut \\
raw_content & texte long & contenu brut \\
collected_at & datetime & date de collecte \\
published_at_detected & datetime & date détectée de publication \\
parsing_status & texte & statut du parsing \\
checksum & texte & empreinte pour détection de doublons \\
created_at & datetime & date de création \\
\hline
\end{tabular}
\end{center}

\subsubsection{Table \texttt{offers}}

Cette table contient les offres normalisées utilisées par le système.

\begin{center}
\begin{tabular}{|l|l|l|}
\hline
\textbf{Champ} & \textbf{Type logique} & \textbf{Description} \\
\hline
id & PK & identifiant interne \\
raw_offer_id & FK & référence vers \texttt{offers\_raw} \\
company_id & FK & référence vers \texttt{companies} \\
primary_source_id & FK & référence vers \texttt{sources} \\
normalized_title & texte & titre normalisé \\
normalized_description & texte long & description nettoyée \\
contract_type & texte & type de contrat \\
duration_months & entier & durée estimée \\
location_text & texte & localisation \\
work_mode & texte & présentiel, hybride, télétravail \\
published_at & datetime & date de publication \\
deadline_at & datetime & date limite éventuelle \\
education_level & texte & niveau requis \\
current_state & texte & état courant de l’offre \\
global_score & numérique & score global \\
action_score & numérique & score d’action \\
score_justification & texte long & justification du scoring \\
is_active & booléen & offre active ou non \\
created_at & datetime & date de création \\
updated_at & datetime & date de mise à jour \\
\hline
\end{tabular}
\end{center}

\subsubsection{Table \texttt{user\_profiles}}

Cette table représente le profil principal utilisé pour l’analyse.

\begin{center}
\begin{tabular}{|l|l|l|}
\hline
\textbf{Champ} & \textbf{Type logique} & \textbf{Description} \\
\hline
id & PK & identifiant interne \\
profile_name & texte & nom du profil \\
education_summary & texte long & résumé de formation \\
skills_text & texte long & compétences \\
experiences_text & texte long & expériences \\
projects_text & texte long & projets \\
constraints_text & texte long & contraintes \\
is_active & booléen & profil actif \\
created_at & datetime & date de création \\
updated_at & datetime & date de mise à jour \\
\hline
\end{tabular}
\end{center}

\subsubsection{Table \texttt{search\_preferences}}

Cette table regroupe les préférences de recherche.

\begin{center}
\begin{tabular}{|l|l|l|}
\hline
\textbf{Champ} & \textbf{Type logique} & \textbf{Description} \\
\hline
id & PK & identifiant interne \\
user_profile_id & FK & référence vers \texttt{user\_profiles} \\
monitoring_frequency_hours & entier & fréquence de surveillance \\
allowed_locations & texte long & localisations autorisées \\
target_companies & texte long & entreprises cibles \\
excluded_companies & texte long & entreprises exclues \\
target_keywords & texte long & mots-clés prioritaires \\
minimum_score_threshold & numérique & seuil minimal \\
followup_policy & texte long & politique de relance \\
created_at & datetime & date de création \\
updated_at & datetime & date de mise à jour \\
\hline
\end{tabular}
\end{center}

\subsubsection{Table \texttt{cv\_versions}}

Cette table stocke les différentes versions de CV.

\begin{center}
\begin{tabular}{|l|l|l|}
\hline
\textbf{Champ} & \textbf{Type logique} & \textbf{Description} \\
\hline
id & PK & identifiant interne \\
user_profile_id & FK & référence vers \texttt{user\_profiles} \\
version_name & texte & nom de la version \\
file_path & texte & chemin ou identifiant fichier \\
orientation_label & texte & orientation métier \\
is_active & booléen & version active \\
created_at & datetime & date de création \\
updated_at & datetime & date de mise à jour \\
\hline
\end{tabular}
\end{center}

\subsubsection{Table \texttt{cover\_letters}}

Cette table stocke les lettres de motivation générées.

\begin{center}
\begin{tabular}{|l|l|l|}
\hline
\textbf{Champ} & \textbf{Type logique} & \textbf{Description} \\
\hline
id & PK & identifiant interne \\
offer_id & FK & référence vers \texttt{offers} \\
content & texte long & contenu de la lettre \\
generation_status & texte & statut de génération \\
validated_at & datetime & date de validation \\
final_version_path & texte & chemin vers version finale \\
created_at & datetime & date de création \\
updated_at & datetime & date de mise à jour \\
\hline
\end{tabular}
\end{center}

\subsubsection{Table \texttt{application\_drafts}}

Cette table représente les brouillons de candidature.

\begin{center}
\begin{tabular}{|l|l|l|}
\hline
\textbf{Champ} & \textbf{Type logique} & \textbf{Description} \\
\hline
id & PK & identifiant interne \\
offer_id & FK & référence vers \texttt{offers} \\
cv_version_id & FK & référence vers \texttt{cv\_versions} \\
cover_letter_id & FK & référence vers \texttt{cover\_letters} \\
prepared_email_subject & texte & sujet préparé \\
prepared_email_body & texte long & corps du message \\
prepared_answers & texte long & réponses formulaires \\
draft_state & texte & état du brouillon \\
created_at & datetime & date de création \\
updated_at & datetime & date de mise à jour \\
\hline
\end{tabular}
\end{center}

\subsubsection{Table \texttt{applications}}

Cette table représente les candidatures effectivement envoyées.

\begin{center}
\begin{tabular}{|l|l|l|}
\hline
\textbf{Champ} & \textbf{Type logique} & \textbf{Description} \\
\hline
id & PK & identifiant interne \\
offer_id & FK & référence vers \texttt{offers} \\
application_draft_id & FK & référence vers \texttt{application\_drafts} \\
submission_channel & texte & email, formulaire, plateforme \\
submitted_at & datetime & date d’envoi \\
application_status & texte & statut courant \\
last_status_change_at & datetime & dernier changement \\
user_notes & texte long & commentaires utilisateur \\
created_at & datetime & date de création \\
updated_at & datetime & date de mise à jour \\
\hline
\end{tabular}
\end{center}

\subsubsection{Table \texttt{followups}}

Cette table stocke les relances préparées ou envoyées.

\begin{center}
\begin{tabular}{|l|l|l|}
\hline
\textbf{Champ} & \textbf{Type logique} & \textbf{Description} \\
\hline
id & PK & identifiant interne \\
application_id & FK & référence vers \texttt{applications} \\
followup_type & texte & type de relance \\
content & texte long & contenu \\
scheduled_at & datetime & date prévue \\
sent_at & datetime & date d’envoi \\
followup_status & texte & statut \\
created_at & datetime & date de création \\
updated_at & datetime & date de mise à jour \\
\hline
\end{tabular}
\end{center}

\subsubsection{Table \texttt{responses}}

Cette table stocke les réponses reçues.

\begin{center}
\begin{tabular}{|l|l|l|}
\hline
\textbf{Champ} & \textbf{Type logique} & \textbf{Description} \\
\hline
id & PK & identifiant interne \\
application_id & FK & référence vers \texttt{applications} \\
channel & texte & canal de réception \\
sender & texte & expéditeur \\
subject & texte & sujet \\
body & texte long & contenu \\
received_at & datetime & date de réception \\
auto_classification & texte & classification automatique \\
created_at & datetime & date de création \\
\hline
\end{tabular}
\end{center}

\subsubsection{Table \texttt{agent\_runs}}

Cette table historise les exécutions du système.

\begin{center}
\begin{tabular}{|l|l|l|}
\hline
\textbf{Champ} & \textbf{Type logique} & \textbf{Description} \\
\hline
id & PK & identifiant interne \\
run_type & texte & type d’exécution \\
started_at & datetime & début d’exécution \\
ended_at & datetime & fin d’exécution \\
detected_offers_count & entier & nombre d’offres détectées \\
retained_offers_count & entier & nombre d’offres retenues \\
run_status & texte & statut final \\
summary_log & texte long & résumé d’exécution \\
created_at & datetime & date de création \\
\hline
\end{tabular}
\end{center}

\subsubsection{Table \texttt{decision\_logs}}

Cette table conserve les décisions du système et de l’utilisateur.

\begin{center}
\begin{tabular}{|l|l|l|}
\hline
\textbf{Champ} & \textbf{Type logique} & \textbf{Description} \\
\hline
id & PK & identifiant interne \\
agent_run_id & FK & référence vers \texttt{agent\_runs} \\
object_type & texte & type d’objet concerné \\
object_id & entier & identifiant de l’objet \\
decision_author & texte & système ou utilisateur \\
decision_type & texte & décision prise \\
decision_reason & texte long & justification \\
metadata_json & texte long & métadonnées associées \\
created_at & datetime & date de décision \\
\hline
\end{tabular}
\end{center}

\subsection{Relations structurantes}

Les relations principales du schéma logique sont les suivantes :

\begin{itemize}
    \item \texttt{sources} 1 --- N \texttt{offers\_raw}
    \item \texttt{companies} 1 --- N \texttt{offers}
    \item \texttt{offers\_raw} 1 --- 0..1 \texttt{offers}
    \item \texttt{sources} 1 --- N \texttt{offers}
    \item \texttt{user\_profiles} 1 --- N \texttt{cv\_versions}
    \item \texttt{user\_profiles} 1 --- N \texttt{search\_preferences}
    \item \texttt{offers} 1 --- 0..1 \texttt{cover\_letters}
    \item \texttt{offers} 1 --- 0..1 \texttt{application\_drafts}
    \item \texttt{application\_drafts} N --- 1 \texttt{cv\_versions}
    \item \texttt{application\_drafts} N --- 0..1 \texttt{cover\_letters}
    \item \texttt{offers} 1 --- 0..1 \texttt{applications}
    \item \texttt{applications} 1 --- N \texttt{followups}
    \item \texttt{applications} 1 --- N \texttt{responses}
    \item \texttt{agent\_runs} 1 --- N \texttt{decision\_logs}
\end{itemize}

\subsection{Contraintes logiques majeures}

Les contraintes logiques suivantes doivent être respectées :

\begin{itemize}
    \item une offre normalisée doit être rattachée à une offre brute d’origine ;
    \item une candidature ne peut être créée que pour une offre existante ;
    \item une relance ne peut exister que pour une candidature déjà soumise ;
    \item toute décision importante doit être historisée dans \texttt{decision\_logs} ;
    \item les doublons d’offres doivent être contrôlés à partir des URLs, identifiants externes et empreintes de contenu ;
    \item une candidature soumise doit conserver une référence vers le brouillon utilisé ;
    \item les fichiers de CV et lettres doivent être versionnés.
\end{itemize}

\subsection{Normalisation future}

Dans une version plus avancée, certaines colonnes textuelles pourront être séparées en tables dédiées afin d’améliorer la normalisation :

\begin{itemize}
    \item table de mots-clés ;
    \item table de compétences ;
    \item table de localisations ;
    \item table de tags d’offres ;
    \item table de canaux de communication ;
    \item table de statuts métier normalisés.
\end{itemize}

\subsection{Remarque de transition}

Le présent schéma relationnel logique constitue une base de travail pour la future implémentation PostgreSQL.

La suite du projet consistera à préciser :

\begin{itemize}
    \item les types SQL exacts ;
    \item les contraintes d’unicité ;
    \item les index ;
    \item les énumérations d’états ;
    \item les stratégies de mise à jour et d’archivage.
\end{itemize}

\section{Workflows et flux opérationnels}

Cette section décrit les principaux workflows métier du système.  
Ces workflows représentent les processus opérationnels exécutés par l’agent personnel de recherche de stage.

\subsection{Workflow de veille des offres}

Le système surveille périodiquement les sources d’offres.

Ce processus est déclenché automatiquement selon une fréquence configurée (par exemple toutes les trois heures).

\begin{enumerate}

\item déclenchement du cycle par le planificateur ;

\item visite des sources actives ;

\item extraction des offres disponibles ;

\item détection des doublons ;

\item normalisation des offres ;

\item stockage dans la base de données.

\end{enumerate}

\subsection{Workflow d'analyse des offres}

Chaque offre collectée est analysée afin d'évaluer sa pertinence.

\begin{enumerate}

\item filtrage initial selon les préférences utilisateur ;

\item extraction des informations importantes :

\begin{itemize}
\item compétences requises ;
\item missions ;
\item localisation ;
\item type de contrat ;
\item durée.
\end{itemize}

\item comparaison avec le profil utilisateur ;

\item calcul du score global ;

\item calcul du score d’action ;

\item décision automatique :

\begin{itemize}
\item ignorer l’offre ;
\item notifier l’utilisateur ;
\item préparer une candidature.
\end{itemize}

\end{enumerate}

\subsection{Workflow de préparation de candidature}

Lorsqu’une offre est jugée pertinente, le système prépare une candidature.

\begin{enumerate}

\item création d’un brouillon de candidature ;

\item sélection de la version de CV la plus appropriée ;

\item génération d’une lettre de motivation personnalisée ;

\item préparation des réponses aux formulaires éventuels ;

\item préparation du message de contact ;

\item changement de statut vers \textbf{A\_VALIDER}.

\end{enumerate}

\subsection{Workflow de validation}

Avant toute soumission, l’utilisateur doit valider la candidature.

Les actions possibles sont :

\begin{itemize}

\item valider et envoyer ;

\item modifier les documents ;

\item rejeter la candidature ;

\item reporter la décision.

\end{itemize}

\subsection{Workflow de soumission}

Une fois validée, la candidature est soumise via le canal approprié.

\begin{itemize}

\item formulaire de plateforme ;

\item email direct ;

\item plateforme de recrutement.

\end{itemize}

La candidature est ensuite enregistrée dans l’état :

\begin{center}
\textbf{EN\_SUIVI}
\end{center}

\subsection{Workflow de suivi}

Après soumission, le système surveille les réponses.

\begin{enumerate}

\item vérification des messages reçus ;

\item association des réponses aux candidatures ;

\item classification automatique :

\begin{itemize}
\item réponse positive ;
\item refus ;
\item demande d'information ;
\item invitation à entretien.
\end{itemize}

\item mise à jour du statut.

\end{enumerate}

\subsection{Workflow de relance}

En absence de réponse dans le délai configuré :

\begin{enumerate}

\item détection de candidature sans réponse ;

\item génération d’un message de relance ;

\item changement de statut vers :

\begin{center}
\textbf{RELANCE\_A\_VALIDER}
\end{center}

\item validation utilisateur ;

\item envoi de la relance ;

\item mise à jour du statut.

\end{enumerate}
\section{Architecture logicielle}

\subsection{Vision générale}

L’architecture logicielle du système repose sur une organisation modulaire en plusieurs couches.

L’objectif est de garantir :

\begin{itemize}
    \item une séparation claire des responsabilités ;
    \item une traçabilité complète des décisions et des actions ;
    \item une évolutivité progressive du système ;
    \item une compatibilité avec une logique de human-in-the-loop ;
    \item une capacité d’orchestration de workflows complexes.
\end{itemize}

Le système est structuré selon les couches suivantes :

\begin{enumerate}
    \item interface utilisateur ;
    \item API applicative ;
    \item services métier ;
    \item orchestration et agents ;
    \item connecteurs externes et automatisation ;
    \item persistance, stockage et observabilité.
\end{enumerate}

\subsection{Schéma logique global}

\begin{center}
\begin{tabular}{c}
Interface utilisateur \\
$\downarrow$ \\
API applicative \\
$\downarrow$ \\
Services métier \\
$\downarrow$ \\
Orchestrateur de workflows et agents \\
$\downarrow$ \\
Connecteurs externes / navigateurs / messagerie \\
$\downarrow$ \\
Base de données / fichiers / logs
\end{tabular}
\end{center}
\subsection{Découpage modulaire recommandé}

Une première structuration logicielle du projet peut être la suivante :

\begin{verbatim}
job_agent/
├── frontend/
│   ├── dashboard/
│   ├── offers/
│   ├── applications/
│   ├── followups/
│   └── settings/
│
├── backend/
│   ├── api/
│   │   ├── routes/
│   │   ├── schemas/
│   │   └── dependencies/
│   │
│   ├── services/
│   │   ├── offer_ingestion_service.py
│   │   ├── matching_service.py
│   │   ├── application_preparation_service.py
│   │   ├── submission_service.py
│   │   ├── followup_service.py
│   │   ├── mailbox_sync_service.py
│   │   └── decision_log_service.py
│   │
│   ├── domain/
│   │   ├── models/
│   │   ├── enums/
│   │   ├── rules/
│   │   └── state_machines/
│   │
│   ├── repositories/
│   │   ├── offer_repository.py
│   │   ├── application_repository.py
│   │   ├── followup_repository.py
│   │   └── decision_log_repository.py
│   │
│   ├── orchestration/
│   │   ├── workflows/
│   │   ├── agents/
│   │   └── schedulers/
│   │
│   ├── connectors/
│   │   ├── linkedin/
│   │   ├── indeed/
│   │   ├── company_pages/
│   │   ├── gmail/
│   │   ├── telegram/
│   │   └── playwright/
│   │
│   └── infrastructure/
│       ├── db/
│       ├── cache/
│       ├── storage/
│       └── logging/
│
└── docs/
\end{verbatim}

\section{Architecture logicielle détaillée}

\subsection{Objectifs de l’architecture}

L’architecture logicielle du système doit répondre aux objectifs suivants :

\begin{itemize}
    \item modularité ;
    \item traçabilité ;
    \item maintenabilité ;
    \item évolutivité ;
    \item compatibilité avec une logique de human-in-the-loop ;
    \item séparation claire entre logique métier, orchestration et exécution externe.
\end{itemize}

\subsection{Découpage architectural}

Le système est structuré en six couches principales :

\begin{enumerate}
    \item couche interface utilisateur ;
    \item couche API applicative ;
    \item couche services métier ;
    \item couche orchestration et agents ;
    \item couche connecteurs externes ;
    \item couche persistance et observabilité.
\end{enumerate}

\subsection{Couche interface utilisateur}

La couche interface utilisateur a pour rôle de permettre à l’utilisateur de :

\begin{itemize}
    \item consulter les offres détectées ;
    \item analyser les scores et justifications ;
    \item visualiser les brouillons de candidature ;
    \item valider ou rejeter les candidatures ;
    \item suivre les statuts des candidatures ;
    \item configurer les préférences du système.
\end{itemize}

Cette couche ne porte pas la logique métier profonde ; elle consomme les services exposés par l’API.

\subsection{Couche API applicative}

La couche API applicative expose les opérations du système.

Elle permet notamment :

\begin{itemize}
    \item la consultation des offres ;
    \item la consultation du pipeline de candidatures ;
    \item le déclenchement de la préparation d’une candidature ;
    \item la validation ou le rejet d’un brouillon ;
    \item la consultation et la modification des paramètres.
\end{itemize}

\subsection{Couche services métier}

La couche services métier contient la logique fonctionnelle centrale du produit.

Elle comprend notamment :

\begin{itemize}
    \item un service d’ingestion et de normalisation des offres ;
    \item un service d’analyse et de scoring ;
    \item un service de préparation de candidature ;
    \item un service de soumission ;
    \item un service de suivi des réponses ;
    \item un service de planification des relances ;
    \item un service de journalisation des décisions.
\end{itemize}

\subsection{Couche orchestration et agents}

La couche d’orchestration coordonne les workflows métier du système.

Elle prend en charge :

\begin{itemize}
    \item les déclenchements planifiés ;
    \item les séquences de traitement multi-étapes ;
    \item les interruptions nécessitant validation humaine ;
    \item la reprise d’exécution après validation ;
    \item l’enchaînement entre collecte, analyse, préparation, suivi et relance.
\end{itemize}

Cette couche doit rester distincte de la logique métier pure.

\subsection{Couche connecteurs externes}

La couche connecteurs externes gère les interactions avec les systèmes extérieurs.

Elle comprend :

\begin{itemize}
    \item les connecteurs vers les plateformes d’offres ;
    \item les connecteurs vers les pages carrières d’entreprises ;
    \item les mécanismes d’automatisation navigateur ;
    \item la connexion à la messagerie ;
    \item les connecteurs de notification.
\end{itemize}

\subsection{Couche persistance et observabilité}

Cette couche assure :

\begin{itemize}
    \item le stockage des objets métier ;
    \item l’historisation des décisions ;
    \item la conservation des documents générés ;
    \item la journalisation des exécutions ;
    \item l’observabilité du système.
\end{itemize}

\subsection{Composants logiciels principaux}

Les composants logiciels principaux identifiés à ce stade sont :

\begin{itemize}
    \item \texttt{Scheduler} ;
    \item \texttt{SourceConnectorManager} ;
    \item \texttt{OfferNormalizer} ;
    \item \texttt{OfferDeduplicator} ;
    \item \texttt{OfferScoringEngine} ;
    \item \texttt{DocumentGenerator} ;
    \item \texttt{ApprovalManager} ;
    \item \texttt{SubmissionExecutor} ;
    \item \texttt{MailboxWatcher} ;
    \item \texttt{FollowUpPlanner} ;
    \item \texttt{AuditLogger}.
\end{itemize}

\subsection{Principes d’implémentation}

Les principes suivants doivent guider l’implémentation :

\begin{itemize}
    \item faible couplage entre modules ;
    \item séparation entre logique métier et infrastructure ;
    \item testabilité unitaire des services ;
    \item centralisation des états métier ;
    \item historisation systématique des décisions importantes ;
    \item capacité de montée en charge progressive.
\end{itemize}

\subsection{Exécution synchrone et asynchrone}

Le système doit combiner des traitements synchrones et asynchrones.

Les opérations synchrones concernent notamment :

\begin{itemize}
    \item la consultation des données ;
    \item la validation utilisateur ;
    \item la navigation dans l’interface.
\end{itemize}

Les opérations asynchrones concernent notamment :

\begin{itemize}
    \item la collecte d’offres ;
    \item la normalisation ;
    \item le scoring ;
    \item la génération de documents ;
    \item la surveillance de la messagerie ;
    \item la planification des relances ;
    \item l’automatisation navigateur.
\end{itemize}

\subsection{Évolution future}

Cette architecture a vocation à être raffinée dans les sections suivantes du projet, notamment :

\begin{itemize}
    \item architecture physique de déploiement ;
    \item choix technologiques détaillés ;
    \item contrats d’API ;
    \item architecture des agents ;
    \item architecture frontend ;
    \item stratégie de sécurité et de confidentialité.
\end{itemize}

\section{Architecture technique concrète}

\subsection{Principes de choix techniques}

L’architecture technique retenue doit permettre :

\begin{itemize}
    \item une implémentation progressive ;
    \item une forte traçabilité ;
    \item une séparation claire entre interface, logique métier, orchestration et exécution externe ;
    \item une compatibilité avec un fonctionnement human-in-the-loop ;
    \item une capacité de reprise des processus longs.
\end{itemize}

Le système sera conçu selon une approche de monolithe modulaire complétée par des workers spécialisés.

\subsection{Stack technique retenue}

La stack technique cible est la suivante :

\begin{itemize}
    \item \textbf{Frontend} : Next.js avec TypeScript ;
    \item \textbf{Backend API} : FastAPI ;
    \item \textbf{Base de données} : PostgreSQL ;
    \item \textbf{Cache et coordination} : Redis ;
    \item \textbf{Orchestration de workflows} : n8n ;
    \item \textbf{Logique agentique durable} : LangGraph ;
    \item \textbf{Automatisation navigateur} : Playwright Python ;
    \item \textbf{Déploiement} : Docker / Docker Compose.
\end{itemize}

\subsection{Rôle des composants}

\subsubsection{Frontend}

Le frontend fournit l’interface utilisateur principale. Il permet :

\begin{itemize}
    \item la consultation des offres ;
    \item la visualisation des scores et justifications ;
    \item la validation ou le rejet des brouillons ;
    \item le suivi des candidatures ;
    \item la configuration du système.
\end{itemize}

\subsubsection{Backend API}

Le backend FastAPI expose les services applicatifs du système.

Il assure :

\begin{itemize}
    \item les endpoints d’accès aux données ;
    \item la validation des entrées ;
    \item le déclenchement d’actions métier ;
    \item la communication entre le frontend et les services internes.
\end{itemize}

\subsubsection{Base de données PostgreSQL}

PostgreSQL stocke les données métier persistantes :

\begin{itemize}
    \item sources ;
    \item offres ;
    \item entreprises ;
    \item brouillons ;
    \item candidatures ;
    \item relances ;
    \item réponses ;
    \item journaux de décision ;
    \item exécutions du système.
\end{itemize}

\subsubsection{Redis}

Redis est utilisé pour :

\begin{itemize}
    \item la coordination légère entre composants ;
    \item le cache ;
    \item la gestion de verrous temporaires ;
    \item le support éventuel de files de tâches.
\end{itemize}

\subsubsection{n8n}

n8n est utilisé comme orchestrateur transverse pour :

\begin{itemize}
    \item les déclenchements planifiés ;
    \item les notifications ;
    \item les workflows simples inter-systèmes ;
    \item certaines pauses nécessitant validation humaine.
\end{itemize}

\subsubsection{LangGraph}

LangGraph est utilisé pour les processus agentiques longs et persistants.

Il prend en charge :

\begin{itemize}
    \item les graphes d’analyse d’offres ;
    \item les graphes de préparation de candidature ;
    \item les interruptions avant soumission ;
    \item la reprise après validation utilisateur ;
    \item la persistance d’état des exécutions longues.
\end{itemize}

\subsubsection{Playwright Python}

Playwright Python est utilisé pour :

\begin{itemize}
    \item l’extraction de données sur des sites dynamiques ;
    \item la navigation automatique sur certaines plateformes ;
    \item le pré-remplissage de formulaires ;
    \item l’assistance à la soumission avant validation finale.
\end{itemize}

\subsection{Architecture de déploiement}

Le système est déployé sous forme de services conteneurisés.

Une première architecture de déploiement comprend :

\begin{itemize}
    \item un conteneur frontend ;
    \item un conteneur backend API ;
    \item un conteneur worker Python ;
    \item un conteneur n8n ;
    \item un conteneur PostgreSQL ;
    \item un conteneur Redis ;
    \item un environnement d’exécution Playwright.
\end{itemize}

\subsection{Mode d’exécution}

Le système combine trois modes d’exécution :

\begin{itemize}
    \item exécution interactive déclenchée par l’utilisateur ;
    \item exécution planifiée déclenchée automatiquement ;
    \item exécution longue avec interruption et reprise.
\end{itemize}

\subsection{Choix d’architecture initiale}

Le projet adopte initialement une architecture de monolithe modulaire.

Ce choix est motivé par :

\begin{itemize}
    \item la simplicité d’implémentation ;
    \item la facilité de maintenance ;
    \item la rapidité de développement ;
    \item l’adéquation avec un produit personnel évolutif.
\end{itemize}

\subsection{Périmètre technique de la première version}

La première version du système comprendra :

\begin{itemize}
    \item une interface web ;
    \item une API applicative ;
    \item une base PostgreSQL ;
    \item un orchestrateur n8n ;
    \item un worker agentique Python ;
    \item un premier ensemble restreint de connecteurs de sources ;
    \item la synchronisation avec la messagerie ;
    \item la notification utilisateur.
\end{itemize}

\section{Plan d’implémentation technique v1}

\subsection{Objectif général}

Le plan d’implémentation technique v1 a pour objectif de construire une première version fonctionnelle du système, couvrant l’ensemble de la chaîne de valeur métier :

\begin{itemize}
    \item collecte des offres ;
    \item normalisation et stockage ;
    \item analyse et scoring ;
    \item préparation des candidatures ;
    \item validation humaine ;
    \item suivi et relance.
\end{itemize}

L’approche retenue est incrémentale.  
Le développement doit suivre un ordre garantissant une montée en valeur progressive du produit.

\subsection{Phase 0 : initialisation du projet}

Cette phase vise à mettre en place le socle technique minimal du projet.

Elle comprend :

\begin{itemize}
    \item la création du dépôt principal ;
    \item la mise en place de la structure des dossiers ;
    \item la configuration de Docker Compose ;
    \item la définition des variables d’environnement ;
    \item l’initialisation du backend FastAPI ;
    \item l’initialisation du frontend Next.js ;
    \item la connexion à PostgreSQL et Redis ;
    \item l’initialisation du système de migrations.
\end{itemize}

\subsection{Phase 1 : noyau domaine et persistance}

Cette phase consiste à implémenter les entités métier et leur persistance.

Elle comprend :

\begin{itemize}
    \item la création des modèles de données ;
    \item l’implémentation des tables principales ;
    \item la définition des états métier ;
    \item la mise en place des repositories ;
    \item la création des premières migrations ;
    \item la création de jeux de données initiaux.
\end{itemize}

\subsection{Phase 2 : ingestion des offres}

Cette phase vise à faire entrer de vraies offres dans le système.

Elle comprend :

\begin{itemize}
    \item l’implémentation d’une interface commune de connecteurs ;
    \item la création d’un gestionnaire de connecteurs ;
    \item le développement des premières sources simples ;
    \item la normalisation des offres collectées ;
    \item la détection initiale des doublons ;
    \item le stockage des offres brutes et normalisées.
\end{itemize}

\subsection{Phase 3 : qualification et scoring}

Cette phase vise à transformer les offres collectées en opportunités évaluées.

Elle comprend :

\begin{itemize}
    \item l’implémentation du moteur de matching ;
    \item l’extraction des informations clés des offres ;
    \item le calcul d’un score global ;
    \item le calcul d’un score d’action ;
    \item la production d’une justification explicite ;
    \item la mise à jour des états d’offre.
\end{itemize}

\subsection{Phase 4 : préparation de candidature}

Cette phase vise à produire un brouillon exploitable pour validation.

Elle comprend :

\begin{itemize}
    \item la sélection d’une version de CV ;
    \item la génération d’une lettre de motivation ;
    \item la préparation d’un message de contact ;
    \item la préparation des réponses à d’éventuelles questions ;
    \item la création des objets de type brouillon de candidature.
\end{itemize}

\subsection{Phase 5 : validation et soumission}

Cette phase vise à intégrer le contrôle humain dans le cycle.

Elle comprend :

\begin{itemize}
    \item la mise en place du mécanisme de validation ;
    \item l’implémentation des transitions d’état associées ;
    \item l’enregistrement des décisions utilisateur ;
    \item la création des candidatures envoyées ;
    \item l’initialisation du pipeline de suivi.
\end{itemize}

\subsection{Phase 6 : suivi et relances}

Cette phase vise à fermer la boucle métier.

Elle comprend :

\begin{itemize}
    \item la synchronisation avec la messagerie ;
    \item la détection des réponses reçues ;
    \item le rattachement des réponses aux candidatures ;
    \item la classification automatique des réponses ;
    \item la planification des relances ;
    \item la génération de brouillons de relance.
\end{itemize}

\subsection{Phase 7 : orchestration}

Cette phase vise à automatiser les déclenchements récurrents du système.

Elle comprend :

\begin{itemize}
    \item la planification périodique de la veille ;
    \item l’envoi de notifications ;
    \item la reprise des workflows après validation ;
    \item l’historisation des exécutions.
\end{itemize}

\subsection{Phase 8 : agentique durable}

Cette phase vise à introduire les premiers graphes d’exécution persistants.

Elle comprend :

\begin{itemize}
    \item l’implémentation de graphes LangGraph ciblés ;
    \item la gestion des interruptions ;
    \item la reprise après validation utilisateur ;
    \item la persistance d’état des processus longs.
\end{itemize}

\subsection{Phase 9 : automatisation navigateur ciblée}

Cette phase vise à introduire l’automatisation de certaines plateformes dynamiques.

Elle comprend :

\begin{itemize}
    \item l’intégration de Playwright Python ;
    \item la création d’un premier connecteur navigateur ;
    \item l’assistance au remplissage de formulaire ;
    \item l’arrêt avant soumission finale.
\end{itemize}

\subsection{Ordre de priorité de développement}

L’ordre de priorité recommandé est le suivant :

\begin{enumerate}
    \item socle technique ;
    \item base de données et domaine ;
    \item ingestion des offres ;
    \item scoring ;
    \item génération de brouillons ;
    \item validation humaine ;
    \item suivi et relances ;
    \item orchestration ;
    \item graphes agentiques ;
    \item automatisation navigateur.
\end{enumerate}

\subsection{Critère de succès de la v1}

La version 1 sera considérée comme satisfaisante si le système est capable de :

\begin{itemize}
    \item surveiller plusieurs sources d’offres ;
    \item enregistrer les nouvelles offres ;
    \item analyser leur pertinence ;
    \item générer un brouillon de candidature ;
    \item soumettre ce brouillon à validation humaine ;
    \item enregistrer une candidature envoyée ;
    \item suivre son statut ;
    \item proposer une relance en absence de réponse.
\end{itemize}

\section{Structuration du dépôt et arborescence technique v1}

\subsection{Principe général}

Le projet est organisé sous la forme d’un mono-dépôt structuré en sous-ensembles cohérents.

Cette organisation vise à :

\begin{itemize}
    \item séparer clairement les responsabilités ;
    \item faciliter le développement incrémental ;
    \item maintenir une bonne lisibilité du code ;
    \item permettre une évolution progressive du système ;
    \item distinguer l’interface utilisateur, l’API, les traitements longs, les workflows et l’infrastructure.
\end{itemize}

\subsection{Organisation générale du dépôt}

Le dépôt principal comprend les sous-ensembles suivants :

\begin{itemize}
    \item \texttt{frontend/} : interface utilisateur ;
    \item \texttt{backend/} : API applicative et logique métier ;
    \item \texttt{worker/} : traitements longs et exécutions agentiques ;
    \item \texttt{workflows/} : workflows d’orchestration ;
    \item \texttt{infra/} : fichiers d’infrastructure et scripts d’environnement ;
    \item \texttt{docs/} : documentation du projet ;
    \item \texttt{shared/} : artefacts partagés ;
    \item \texttt{data/} : jeux de données, fixtures et exemples.
\end{itemize}

\subsection{Arborescence cible du dépôt}

Une arborescence technique cible pour la version 1 est la suivante :

\begin{verbatim}
job-search-agent/
├── docs/
├── infra/
├── frontend/
├── backend/
├── worker/
├── workflows/
├── shared/
└── data/
\end{verbatim}

\subsection{Structuration du backend}

Le backend est organisé autour des modules suivants :

\begin{itemize}
    \item \texttt{core/} : configuration, base de données, logging ;
    \item \texttt{api/} : routes et dépendances HTTP ;
    \item \texttt{schemas/} : schémas de validation et de sérialisation ;
    \item \texttt{domain/} : entités, règles métier, états ;
    \item \texttt{services/} : logique applicative ;
    \item \texttt{repositories/} : accès aux données ;
    \item \texttt{orchestration/} : commandes, événements et gestion des flux ;
    \item \texttt{connectors/} : intégrations avec les systèmes externes ;
    \item \texttt{infrastructure/} : implémentations techniques concrètes.
\end{itemize}

\subsection{Structuration du worker}

Le worker regroupe les traitements longs et les processus non interactifs.

Il comprend notamment :

\begin{itemize}
    \item des jobs planifiés ;
    \item des graphes d’exécution ;
    \item des tâches unitaires ;
    \item des intégrations avec les autres composants ;
    \item des utilitaires d’exécution.
\end{itemize}

\subsection{Structuration du frontend}

Le frontend est structuré par domaines fonctionnels.

Il comprend notamment :

\begin{itemize}
    \item les pages principales de l’application ;
    \item les composants d’interface réutilisables ;
    \item les modules fonctionnels par feature ;
    \item les hooks et bibliothèques utilitaires ;
    \item les types frontaux.
\end{itemize}

\subsection{Structuration des workflows}

Les workflows d’orchestration sont stockés séparément afin de permettre :

\begin{itemize}
    \item leur versionnement ;
    \item leur documentation ;
    \item leur export et leur réimportation ;
    \item leur maintenance indépendante.
\end{itemize}

\subsection{Structuration de la documentation}

La documentation doit être maintenue au sein du dépôt.

Elle comprend notamment :

\begin{itemize}
    \item le cahier des charges ;
    \item les décisions d’architecture ;
    \item les notes de workflows ;
    \item les diagrammes ;
    \item les conventions de développement.
\end{itemize}

\subsection{Fichiers racine indispensables}

Les fichiers racine à prévoir dès la version 1 sont :

\begin{itemize}
    \item \texttt{README.md} ;
    \item \texttt{.env.example} ;
    \item \texttt{docker-compose.yml} ;
    \item \texttt{Makefile}.
\end{itemize}

\subsection{Principes de structuration}

Les principes suivants doivent être respectés :

\begin{itemize}
    \item un module doit avoir une responsabilité claire ;
    \item la logique métier doit être séparée de l’infrastructure ;
    \item les intégrations externes doivent être isolées ;
    \item les états métier doivent être centralisés ;
    \item les artefacts partagés doivent être explicitement identifiés.
\end{itemize}

\section{Catalogue des modules backend v1}

\subsection{Objectif}

Cette section recense les principaux modules backend de la version 1 du système.

Pour chaque module, l’objectif est de préciser :

\begin{itemize}
    \item sa responsabilité ;
    \item ses entrées ;
    \item ses sorties ;
    \item ses dépendances principales ;
    \item sa priorité d’implémentation.
\end{itemize}

\subsection{Familles de modules}

Les modules backend v1 sont regroupés en six familles :

\begin{enumerate}
    \item modules cœur technique ;
    \item modules de persistance ;
    \item modules métier liés aux offres ;
    \item modules métier liés aux candidatures ;
    \item modules métier liés au suivi ;
    \item modules transverses.
\end{enumerate}

\subsection{Modules cœur technique}

\subsubsection{\texttt{config}}
Responsabilité : centraliser la configuration applicative.

\subsubsection{\texttt{database}}
Responsabilité : gérer l’accès à la base de données et les sessions.

\subsubsection{\texttt{logging}}
Responsabilité : centraliser la journalisation applicative.

\subsubsection{\texttt{enums}}
Responsabilité : définir les états et catégories métier.

\subsubsection{\texttt{state\_machines}}
Responsabilité : encapsuler les transitions légales d’état.

\subsection{Modules de persistance}

\subsubsection{\texttt{source\_repository}}
Responsabilité : gérer les sources surveillées.

\subsubsection{\texttt{company\_repository}}
Responsabilité : gérer les entreprises.

\subsubsection{\texttt{offer\_raw\_repository}}
Responsabilité : gérer les offres brutes collectées.

\subsubsection{\texttt{offer\_repository}}
Responsabilité : gérer les offres normalisées.

\subsubsection{\texttt{application\_draft\_repository}}
Responsabilité : gérer les brouillons de candidature.

\subsubsection{\texttt{application\_repository}}
Responsabilité : gérer les candidatures envoyées.

\subsubsection{\texttt{followup\_repository}}
Responsabilité : gérer les relances.

\subsubsection{\texttt{response\_repository}}
Responsabilité : gérer les réponses reçues.

\subsubsection{\texttt{decision\_log\_repository}}
Responsabilité : gérer la persistance des journaux de décision.

\subsection{Modules métier liés aux offres}

\subsubsection{\texttt{source\_connector\_manager}}
Responsabilité : piloter les connecteurs de collecte.

\subsubsection{\texttt{offer\_normalizer}}
Responsabilité : transformer une offre brute en offre normalisée.

\subsubsection{\texttt{offer\_deduplicator}}
Responsabilité : détecter les doublons d’offres.

\subsubsection{\texttt{offer\_ingestion\_service}}
Responsabilité : orchestrer l’entrée d’une offre dans le système.

\subsubsection{\texttt{matching\_service}}
Responsabilité : évaluer la compatibilité entre une offre et le profil utilisateur.

\subsubsection{\texttt{offer\_scoring\_engine}}
Responsabilité : calculer les scores métier.

\subsection{Modules métier liés aux candidatures}

\subsubsection{\texttt{cv\_selection\_service}}
Responsabilité : choisir la version de CV la plus adaptée.

\subsubsection{\texttt{cover\_letter\_service}}
Responsabilité : générer une lettre de motivation adaptée.

\subsubsection{\texttt{application\_preparation\_service}}
Responsabilité : préparer un brouillon complet de candidature.

\subsubsection{\texttt{approval\_service}}
Responsabilité : gérer les validations utilisateur.

\subsubsection{\texttt{submission\_service}}
Responsabilité : enregistrer ou exécuter la soumission d’une candidature validée.

\subsection{Modules métier liés au suivi}

\subsubsection{\texttt{mailbox\_sync\_service}}
Responsabilité : synchroniser les réponses issues de la messagerie.

\subsubsection{\texttt{response\_classification\_service}}
Responsabilité : classifier automatiquement les réponses reçues.

\subsubsection{\texttt{followup\_policy\_service}}
Responsabilité : déterminer si une relance est autorisée ou recommandée.

\subsubsection{\texttt{followup\_service}}
Responsabilité : créer et gérer les brouillons de relance.

\subsection{Modules transverses}

\subsubsection{\texttt{decision\_log\_service}}
Responsabilité : centraliser la journalisation métier.

\subsubsection{\texttt{agent\_run\_service}}
Responsabilité : historiser les exécutions du système.

\subsubsection{\texttt{notification\_service}}
Responsabilité : envoyer les notifications utiles à l’utilisateur.

\subsection{Ordre d’implémentation recommandé}

L’ordre d’implémentation recommandé est le suivant :

\begin{enumerate}
    \item modules cœur technique ;
    \item modules de persistance ;
    \item modules d’ingestion et de normalisation des offres ;
    \item modules de scoring et de matching ;
    \item modules de préparation de candidature ;
    \item modules de validation et de soumission ;
    \item modules de suivi et de relance ;
    \item modules transverses avancés.
\end{enumerate}

\subsection{Modules critiques de la version 1}

Les modules critiques de la version 1 sont :

\begin{itemize}
    \item \texttt{offer\_ingestion\_service} ;
    \item \texttt{offer\_normalizer} ;
    \item \texttt{offer\_deduplicator} ;
    \item \texttt{matching\_service} ;
    \item \texttt{application\_preparation\_service} ;
    \item \texttt{approval\_service} ;
    \item \texttt{submission\_service}.
\end{itemize}

\section{Spécification des interfaces et contrats des modules critiques v1}

\subsection{Objectif}

Cette section précise les contrats logiciels des modules critiques de la version 1.

Pour chaque module, il convient de définir :

\begin{itemize}
    \item sa responsabilité ;
    \item ses préconditions ;
    \item ses entrées ;
    \item ses sorties ;
    \item ses postconditions ;
    \item ses erreurs possibles ;
    \item ses dépendances principales.
\end{itemize}

\subsection{\texttt{BaseSourceConnector}}

\textbf{Responsabilité :} définir le contrat minimal de tout connecteur de source.

\textbf{Entrées :}
\begin{itemize}
    \item configuration de source ;
    \item paramètres éventuels de collecte.
\end{itemize}

\textbf{Sortie :}
\begin{itemize}
    \item liste d’offres brutes standardisées.
\end{itemize}

\textbf{Garantie :}
chaque offre brute doit contenir un noyau minimal d’information permettant son ingestion.

\subsection{\texttt{OfferNormalizer}}

\textbf{Responsabilité :} transformer une offre brute en offre normalisée.

\textbf{Entrée :}
\begin{itemize}
    \item une offre brute.
\end{itemize}

\textbf{Sortie :}
\begin{itemize}
    \item une offre normalisée structurée.
\end{itemize}

\textbf{Postconditions :}
la sortie doit contenir autant que possible un titre normalisé, une entreprise, une description nettoyée, une localisation, un type de contrat, une date de publication et les métadonnées utiles.

\subsection{\texttt{OfferDeduplicator}}

\textbf{Responsabilité :} identifier si une offre correspond à un nouvel enregistrement, à un doublon ou à une mise à jour.

\textbf{Entrées :}
\begin{itemize}
    \item une offre normalisée ;
    \item le contexte de comparaison en base.
\end{itemize}

\textbf{Sortie :}
\begin{itemize}
    \item une décision de déduplication ;
    \item un score de confiance ;
    \item éventuellement une référence vers l’offre existante.
\end{itemize}

\subsection{\texttt{OfferIngestionService}}

\textbf{Responsabilité :} orchestrer l’entrée complète d’une offre dans le système.

\textbf{Entrée :}
\begin{itemize}
    \item une offre brute ;
    \item une référence de source.
\end{itemize}

\textbf{Sortie :}
\begin{itemize}
    \item un résultat d’ingestion contenant les identifiants créés ou mis à jour ;
    \item le statut d’ingestion ;
    \item la décision de déduplication.
\end{itemize}

\textbf{Postconditions :}
l’offre brute est conservée, l’offre normalisée est créée ou mise à jour, et la décision est historisée.

\subsection{\texttt{MatchingService}}

\textbf{Responsabilité :} évaluer la pertinence d’une offre relativement au profil utilisateur.

\textbf{Entrées :}
\begin{itemize}
    \item une offre normalisée ;
    \item un profil utilisateur ;
    \item les préférences de recherche actives.
\end{itemize}

\textbf{Sortie :}
\begin{itemize}
    \item un résultat de matching contenant les sous-scores ;
    \item un score global ;
    \item un score d’action ;
    \item une justification ;
    \item une recommandation d’action.
\end{itemize}

\textbf{Garantie :}
le résultat doit être explicable, traçable et réutilisable dans les décisions métier.

\subsection{\texttt{ApplicationPreparationService}}

\textbf{Responsabilité :} préparer un brouillon complet de candidature.

\textbf{Entrées :}
\begin{itemize}
    \item une offre qualifiée ;
    \item un profil utilisateur ;
    \item des préférences actives ;
    \item les versions de CV disponibles.
\end{itemize}

\textbf{Sortie :}
\begin{itemize}
    \item un brouillon de candidature complet ;
    \item la référence du CV sélectionné ;
    \item la lettre générée ;
    \item les messages préparés ;
    \item l’état \texttt{A\_VALIDER}.
\end{itemize}

\textbf{Postconditions :}
un brouillon persistant est créé, les artefacts sont stockés et la décision est journalisée.

\subsection{\texttt{ApprovalService}}

\textbf{Responsabilité :} traiter les décisions utilisateur sur les objets à valider.

\textbf{Entrées :}
\begin{itemize}
    \item type d’objet ;
    \item identifiant d’objet ;
    \item décision utilisateur ;
    \item commentaire éventuel.
\end{itemize}

\textbf{Sortie :}
\begin{itemize}
    \item résultat de validation ;
    \item ancien état ;
    \item nouvel état ;
    \item journal de décision.
\end{itemize}

\textbf{Garantie :}
aucune transition illégale ne doit être autorisée.

\subsection{\texttt{SubmissionService}}

\textbf{Responsabilité :} transformer un brouillon validé en candidature soumise.

\textbf{Entrées :}
\begin{itemize}
    \item un brouillon validé ;
    \item un canal de soumission ;
    \item les métadonnées éventuelles.
\end{itemize}

\textbf{Sortie :}
\begin{itemize}
    \item une candidature créée ;
    \item un statut initial de suivi ;
    \item une trace de soumission.
\end{itemize}

\textbf{Postconditions :}
la candidature est créée, le brouillon est marqué comme soumis et le pipeline de suivi est initialisé.

\subsection{Structures contractuelles minimales}

Les structures minimales à formaliser sont notamment :

\begin{itemize}
    \item \texttt{RawOfferPayload} ;
    \item \texttt{NormalizedOfferPayload} ;
    \item \texttt{DeduplicationDecision} ;
    \item \texttt{MatchingResult} ;
    \item \texttt{ApplicationDraftResult} ;
    \item \texttt{ApprovalResult} ;
    \item \texttt{SubmissionResult}.
\end{itemize}

\subsection{Règles générales}

Les règles suivantes s’appliquent à tous les modules critiques :

\begin{itemize}
    \item une responsabilité fonctionnelle claire par module ;
    \item aucune transition d’état hors machine d’état ;
    \item journalisation systématique des décisions importantes ;
    \item sorties structurées et explicites.
\end{itemize}

\section{Machines d’état métier v1}

\subsection{Objectif}

Cette section formalise les machines d’état métier de la version 1 du système.

L’objectif est de garantir :

\begin{itemize}
    \item la cohérence des cycles de vie des objets métier ;
    \item le contrôle des transitions autorisées ;
    \item la traçabilité des changements d’état ;
    \item l’intégration correcte du human-in-the-loop.
\end{itemize}

Les objets métier principaux soumis à une machine d’état sont :

\begin{itemize}
    \item les offres ;
    \item les brouillons de candidature ;
    \item les candidatures ;
    \item les relances.
\end{itemize}

\subsection{Principes généraux}

Les principes suivants s’appliquent à toutes les machines d’état :

\begin{itemize}
    \item toute transition d’état doit passer par la machine d’état correspondante ;
    \item toute transition importante doit être journalisée ;
    \item un état doit représenter une réalité métier explicite ;
    \item les états terminaux ne peuvent pas être quittés dans la version 1.
\end{itemize}

\subsection{Machine d’état des offres}

Les états métier d’une offre sont les suivants :

\begin{itemize}
    \item \texttt{DETECTED} ;
    \item \texttt{RAW\_STORED} ;
    \item \texttt{NORMALIZED} ;
    \item \texttt{DEDUPLICATED} ;
    \item \texttt{ANALYZED} ;
    \item \texttt{REJECTED} ;
    \item \texttt{QUALIFIED} ;
    \item \texttt{DRAFT\_REQUESTED} ;
    \item \texttt{DRAFT\_PREPARED} ;
    \item \texttt{READY\_FOR\_REVIEW} ;
    \item \texttt{SUBMISSION\_IN\_PROGRESS} ;
    \item \texttt{SUBMITTED} ;
    \item \texttt{CLOSED}.
\end{itemize}

Le flux principal est :

\begin{center}
\texttt{DETECTED}
$\rightarrow$
\texttt{RAW\_STORED}
$\rightarrow$
\texttt{NORMALIZED}
$\rightarrow$
\texttt{DEDUPLICATED}
$\rightarrow$
\texttt{ANALYZED}
$\rightarrow$
\texttt{QUALIFIED}
$\rightarrow$
\texttt{DRAFT\_REQUESTED}
$\rightarrow$
\texttt{DRAFT\_PREPARED}
$\rightarrow$
\texttt{READY\_FOR\_REVIEW}
$\rightarrow$
\texttt{SUBMISSION\_IN\_PROGRESS}
$\rightarrow$
\texttt{SUBMITTED}
$\rightarrow$
\texttt{CLOSED}
\end{center}

Des branches alternatives sont autorisées :

\begin{itemize}
    \item \texttt{ANALYZED} $\rightarrow$ \texttt{REJECTED} ;
    \item \texttt{QUALIFIED} $\rightarrow$ \texttt{REJECTED} ;
    \item \texttt{READY\_FOR\_REVIEW} $\rightarrow$ \texttt{REJECTED}.
\end{itemize}

\subsection{Machine d’état des brouillons de candidature}

Les états métier d’un brouillon sont les suivants :

\begin{itemize}
    \item \texttt{DRAFT\_CREATED} ;
    \item \texttt{CV\_SELECTED} ;
    \item \texttt{LETTER\_GENERATED} ;
    \item \texttt{MESSAGE\_PREPARED} ;
    \item \texttt{READY\_FOR\_VALIDATION} ;
    \item \texttt{APPROVED} ;
    \item \texttt{REJECTED} ;
    \item \texttt{CHANGES\_REQUESTED} ;
    \item \texttt{POSTPONED} ;
    \item \texttt{SUBMITTED}.
\end{itemize}

Le flux principal est :

\begin{center}
\texttt{DRAFT\_CREATED}
$\rightarrow$
\texttt{CV\_SELECTED}
$\rightarrow$
\texttt{LETTER\_GENERATED}
$\rightarrow$
\texttt{MESSAGE\_PREPARED}
$\rightarrow$
\texttt{READY\_FOR\_VALIDATION}
$\rightarrow$
\texttt{APPROVED}
$\rightarrow$
\texttt{SUBMITTED}
\end{center}

Des branches alternatives sont autorisées :

\begin{itemize}
    \item \texttt{READY\_FOR\_VALIDATION} $\rightarrow$ \texttt{REJECTED} ;
    \item \texttt{READY\_FOR\_VALIDATION} $\rightarrow$ \texttt{CHANGES\_REQUESTED} ;
    \item \texttt{READY\_FOR\_VALIDATION} $\rightarrow$ \texttt{POSTPONED} ;
    \item \texttt{CHANGES\_REQUESTED} $\rightarrow$ \texttt{DRAFT\_CREATED} ;
    \item \texttt{POSTPONED} $\rightarrow$ \texttt{READY\_FOR\_VALIDATION}.
\end{itemize}

\subsection{Machine d’état des candidatures}

Les états métier d’une candidature sont les suivants :

\begin{itemize}
    \item \texttt{CREATED} ;
    \item \texttt{SUBMITTED} ;
    \item \texttt{ACKNOWLEDGED} ;
    \item \texttt{UNDER\_REVIEW} ;
    \item \texttt{FOLLOWUP\_DUE} ;
    \item \texttt{FOLLOWUP\_IN\_PROGRESS} ;
    \item \texttt{RESPONSE\_RECEIVED} ;
    \item \texttt{INTERVIEW} ;
    \item \texttt{REJECTED} ;
    \item \texttt{ACCEPTED} ;
    \item \texttt{WITHDRAWN} ;
    \item \texttt{CLOSED}.
\end{itemize}

Des transitions principales sont autorisées entre :

\begin{itemize}
    \item \texttt{CREATED} $\rightarrow$ \texttt{SUBMITTED} ;
    \item \texttt{SUBMITTED} $\rightarrow$ \texttt{ACKNOWLEDGED} ;
    \item \texttt{SUBMITTED} $\rightarrow$ \texttt{FOLLOWUP\_DUE} ;
    \item \texttt{SUBMITTED} $\rightarrow$ \texttt{REJECTED} ;
    \item \texttt{SUBMITTED} $\rightarrow$ \texttt{ACCEPTED} ;
    \item \texttt{ACKNOWLEDGED} $\rightarrow$ \texttt{UNDER\_REVIEW} ;
    \item \texttt{UNDER\_REVIEW} $\rightarrow$ \texttt{RESPONSE\_RECEIVED} ;
    \item \texttt{RESPONSE\_RECEIVED} $\rightarrow$ \texttt{INTERVIEW} ;
    \item \texttt{INTERVIEW} $\rightarrow$ \texttt{REJECTED} ;
    \item \texttt{INTERVIEW} $\rightarrow$ \texttt{ACCEPTED} ;
    \item \texttt{REJECTED} $\rightarrow$ \texttt{CLOSED} ;
    \item \texttt{ACCEPTED} $\rightarrow$ \texttt{CLOSED} ;
    \item \texttt{WITHDRAWN} $\rightarrow$ \texttt{CLOSED}.
\end{itemize}

\subsection{Machine d’état des relances}

Les états métier d’une relance sont les suivants :

\begin{itemize}
    \item \texttt{NOT\_NEEDED} ;
    \item \texttt{ELIGIBLE} ;
    \item \texttt{DRAFTING} ;
    \item \texttt{READY\_FOR\_VALIDATION} ;
    \item \texttt{APPROVED} ;
    \item \texttt{SENT} ;
    \item \texttt{CANCELLED} ;
    \item \texttt{ANSWERED} ;
    \item \texttt{CLOSED}.
\end{itemize}

Le flux principal est :

\begin{center}
\texttt{NOT\_NEEDED}
$\rightarrow$
\texttt{ELIGIBLE}
$\rightarrow$
\texttt{DRAFTING}
$\rightarrow$
\texttt{READY\_FOR\_VALIDATION}
$\rightarrow$
\texttt{APPROVED}
$\rightarrow$
\texttt{SENT}
$\rightarrow$
\texttt{CLOSED}
\end{center}

Des branches alternatives sont autorisées :

\begin{itemize}
    \item \texttt{READY\_FOR\_VALIDATION} $\rightarrow$ \texttt{CANCELLED} ;
    \item \texttt{SENT} $\rightarrow$ \texttt{ANSWERED} ;
    \item \texttt{ANSWERED} $\rightarrow$ \texttt{CLOSED} ;
    \item \texttt{CANCELLED} $\rightarrow$ \texttt{CLOSED}.
\end{itemize}

\subsection{Règles d’interdiction}

Les interdictions suivantes sont fixées pour la version 1 :

\begin{itemize}
    \item une offre ne peut pas passer directement de \texttt{DETECTED} à \texttt{SUBMITTED} ;
    \item un brouillon ne peut pas être soumis sans être approuvé ;
    \item une candidature rejetée, acceptée ou clôturée ne peut pas générer de nouvelle relance ;
    \item une relance annulée ou clôturée ne peut pas être réactivée.
\end{itemize}

\subsection{Événements métier}

Les transitions sont déclenchées par des événements métier explicites, tels que :

\begin{itemize}
    \item détection d’offre ;
    \item normalisation terminée ;
    \item scoring terminé ;
    \item brouillon préparé ;
    \item validation utilisateur ;
    \item soumission effectuée ;
    \item réponse reçue ;
    \item relance devenue éligible ;
    \item relance envoyée.
\end{itemize}

\subsection{Règle générale de journalisation}

Toute transition importante d’état doit être historisée dans les journaux de décision ou d’exécution du système.

\section{Spécification des DTO, payloads et schémas d’échange v1}

\subsection{Objectif}

Cette section formalise les principaux objets d’échange de la version 1 du système.

Ces objets ont pour rôle de structurer les échanges :

\begin{itemize}
    \item entre connecteurs et backend ;
    \item entre services métier ;
    \item entre backend et worker ;
    \item entre backend et frontend ;
    \item entre le système et les services externes.
\end{itemize}

\subsection{Principes généraux}

Les principes suivants s’appliquent aux DTO de la version 1 :

\begin{itemize}
    \item un DTO représente un contrat d’échange explicite ;
    \item un DTO doit être simple, sérialisable et stable ;
    \item les champs critiques doivent être obligatoires ;
    \item les états métier doivent être portés par des énumérations ;
    \item un champ de métadonnées extensible peut être ajouté si nécessaire.
\end{itemize}

\subsection{DTO d’ingestion des offres}

Les objets principaux de cette famille sont :

\begin{itemize}
    \item \texttt{SourceConfigDTO} ;
    \item \texttt{RawOfferPayload} ;
    \item \texttt{ConnectorRunResult}.
\end{itemize}

\texttt{SourceConfigDTO} décrit une source surveillée.

\texttt{RawOfferPayload} représente une offre brute retournée par un connecteur.

\texttt{ConnectorRunResult} représente le résultat d’un cycle d’exécution d’un connecteur.

\subsection{DTO de normalisation et de déduplication}

Les objets principaux de cette famille sont :

\begin{itemize}
    \item \texttt{NormalizedOfferPayload} ;
    \item \texttt{NormalizationResult} ;
    \item \texttt{DeduplicationDecision} ;
    \item \texttt{OfferIngestionResult}.
\end{itemize}

\texttt{NormalizedOfferPayload} représente une offre après structuration.

\texttt{NormalizationResult} décrit le résultat du processus de normalisation.

\texttt{DeduplicationDecision} formalise la décision de déduplication.

\texttt{OfferIngestionResult} formalise le résultat complet d’une ingestion.

\subsection{DTO de matching et de scoring}

Les objets principaux de cette famille sont :

\begin{itemize}
    \item \texttt{MatchingInputDTO} ;
    \item \texttt{DimensionScoreDTO} ;
    \item \texttt{MatchingResult} ;
    \item \texttt{OfferScoringSnapshotDTO}.
\end{itemize}

\texttt{MatchingResult} contient notamment :

\begin{itemize}
    \item des sous-scores par dimension ;
    \item un score global ;
    \item un score d’action ;
    \item une justification ;
    \item une recommandation d’action.
\end{itemize}

\subsection{DTO de préparation et de validation de candidature}

Les objets principaux de cette famille sont :

\begin{itemize}
    \item \texttt{CVSelectionResult} ;
    \item \texttt{CoverLetterDraftDTO} ;
    \item \texttt{PreparedMessageDTO} ;
    \item \texttt{PreparedAnswersDTO} ;
    \item \texttt{ApplicationDraftResult} ;
    \item \texttt{ApprovalRequestDTO} ;
    \item \texttt{ApprovalResult} ;
    \item \texttt{SubmissionRequestDTO} ;
    \item \texttt{SubmissionResult}.
\end{itemize}

\texttt{ApplicationDraftResult} représente le résultat complet de préparation d’un brouillon.

\texttt{ApprovalRequestDTO} et \texttt{ApprovalResult} structurent les décisions utilisateur.

\texttt{SubmissionRequestDTO} et \texttt{SubmissionResult} structurent la soumission d’une candidature.

\subsection{DTO de suivi, réponses et relances}

Les objets principaux de cette famille sont :

\begin{itemize}
    \item \texttt{MailboxMessageDTO} ;
    \item \texttt{ResponseClassificationResult} ;
    \item \texttt{FollowUpEligibilityResult} ;
    \item \texttt{FollowUpDraftDTO} ;
    \item \texttt{FollowUpResult}.
\end{itemize}

Ces objets permettent de formaliser :

\begin{itemize}
    \item les messages reçus ;
    \item leur classification ;
    \item l’éligibilité à une relance ;
    \item les brouillons de relance.
\end{itemize}

\subsection{DTO transverses}

Les objets transverses principaux sont :

\begin{itemize}
    \item \texttt{DecisionLogDTO} ;
    \item \texttt{AgentRunDTO} ;
    \item \texttt{NotificationDTO}.
\end{itemize}

Ils permettent respectivement de représenter :

\begin{itemize}
    \item une décision métier historisée ;
    \item une exécution du système ;
    \item une notification à destination de l’utilisateur.
\end{itemize}

\subsection{Règles de conception}

Les règles suivantes doivent être respectées :

\begin{itemize}
    \item les DTO internes, les DTO d’API et les DTO d’intégration externe doivent être distingués ;
    \item tout objet traversant plusieurs modules doit être formalisé explicitement ;
    \item les payloads doivent rester sobres dans la version 1 ;
    \item les sorties critiques doivent être fortement structurées.
\end{itemize}

\subsection{Priorités pour la version 1}

Les DTO prioritaires à stabiliser pour la version 1 sont notamment :

\begin{itemize}
    \item \texttt{RawOfferPayload} ;
    \item \texttt{NormalizedOfferPayload} ;
    \item \texttt{DeduplicationDecision} ;
    \item \texttt{MatchingResult} ;
    \item \texttt{ApplicationDraftResult} ;
    \item \texttt{ApprovalRequestDTO} ;
    \item \texttt{ApprovalResult} ;
    \item \texttt{SubmissionRequestDTO} ;
    \item \texttt{SubmissionResult} ;
    \item \texttt{FollowUpEligibilityResult}.
\end{itemize}

\section{Spécification des endpoints API v1}

\subsection{Objectif}

Cette section décrit la surface d’API de la version 1 du système.

L’objectif de cette API est de permettre :

\begin{itemize}
    \item la consultation des offres détectées ;
    \item la consultation détaillée des objets métier ;
    \item la préparation de candidatures ;
    \item la validation humaine ;
    \item la soumission et le suivi des candidatures ;
    \item la préparation et la gestion des relances ;
    \item la consultation des préférences et des exécutions du système.
\end{itemize}

\subsection{Principes généraux}

L’API de la version 1 suit les principes suivants :

\begin{itemize}
    \item style REST simple ;
    \item endpoints orientés ressources métier ;
    \item cohérence avec les machines d’état ;
    \item séparation entre opérations de lecture et opérations de transition d’état ;
    \item réponses structurées et traçables.
\end{itemize}

\subsection{Familles d’endpoints}

Les familles principales d’endpoints sont les suivantes :

\begin{itemize}
    \item \texttt{/health} ;
    \item \texttt{/offers} ;
    \item \texttt{/drafts} ;
    \item \texttt{/applications} ;
    \item \texttt{/followups} ;
    \item \texttt{/profiles} ;
    \item \texttt{/settings} ;
    \item \texttt{/runs}.
\end{itemize}

\subsection{Endpoints de santé}

\subsubsection{\texttt{GET /health}}
Permet de vérifier que l’API répond correctement.

\subsubsection{\texttt{GET /health/dependencies}}
Permet de vérifier l’état des dépendances principales.

\subsection{Endpoints liés aux offres}

\subsubsection{\texttt{GET /offers}}
Liste les offres avec filtres, tri et pagination.

\subsubsection{\texttt{GET /offers/\{offer\_id\}}}
Retourne le détail d’une offre.

\subsubsection{\texttt{POST /offers/\{offer\_id\}/score}}
Déclenche ou redéclenche le scoring d’une offre.

\subsubsection{\texttt{POST /offers/\{offer\_id\}/prepare-draft}}
Demande la préparation d’un brouillon de candidature.

\subsubsection{\texttt{POST /offers/\{offer\_id\}/reject}}
Permet de rejeter explicitement une offre.

\subsubsection{\texttt{GET /offers/\{offer\_id\}/history}}
Retourne l’historique métier d’une offre.

\subsection{Endpoints liés aux brouillons de candidature}

\subsubsection{\texttt{GET /drafts}}
Liste les brouillons de candidature.

\subsubsection{\texttt{GET /drafts/\{draft\_id\}}}
Retourne le détail d’un brouillon.

\subsubsection{\texttt{POST /drafts/\{draft\_id\}/approve}}
Approuve un brouillon.

\subsubsection{\texttt{POST /drafts/\{draft\_id\}/reject}}
Rejette un brouillon.

\subsubsection{\texttt{POST /drafts/\{draft\_id\}/request-changes}}
Demande une modification d’un brouillon.

\subsubsection{\texttt{POST /drafts/\{draft\_id\}/postpone}}
Reporte la décision sur un brouillon.

\subsubsection{\texttt{POST /drafts/\{draft\_id\}/submit}}
Transforme un brouillon validé en candidature suivie.

\subsection{Endpoints liés aux candidatures}

\subsubsection{\texttt{GET /applications}}
Liste les candidatures et leur état de suivi.

\subsubsection{\texttt{GET /applications/\{application\_id\}}}
Retourne le détail d’une candidature.

\subsubsection{\texttt{POST /applications/\{application\_id\}/mark-status}}
Permet de confirmer ou corriger un changement de statut.

\subsubsection{\texttt{GET /applications/\{application\_id\}/history}}
Retourne l’historique complet d’une candidature.

\subsection{Endpoints liés aux relances}

\subsubsection{\texttt{GET /followups}}
Liste les relances.

\subsubsection{\texttt{GET /applications/\{application\_id\}/followup-eligibility}}
Évalue l’éligibilité d’une candidature à une relance.

\subsubsection{\texttt{POST /applications/\{application\_id\}/followups/prepare}}
Prépare une relance.

\subsubsection{\texttt{GET /followups/\{followup\_id\}}}
Retourne le détail d’une relance.

\subsubsection{\texttt{POST /followups/\{followup\_id\}/approve}}
Approuve une relance.

\subsubsection{\texttt{POST /followups/\{followup\_id\}/cancel}}
Annule une relance.

\subsubsection{\texttt{POST /followups/\{followup\_id\}/send}}
Déclenche ou enregistre l’envoi d’une relance.

\subsection{Endpoints liés au profil utilisateur}

\subsubsection{\texttt{GET /profiles/active}}
Retourne le profil utilisateur actif.

\subsubsection{\texttt{PUT /profiles/active}}
Met à jour le profil utilisateur actif.

\subsubsection{\texttt{GET /profiles/active/cv-versions}}
Liste les versions de CV disponibles.

\subsubsection{\texttt{POST /profiles/active/cv-versions}}
Ajoute une nouvelle version de CV.

\subsection{Endpoints liés aux préférences}

\subsubsection{\texttt{GET /settings/search-preferences}}
Retourne les préférences de recherche.

\subsubsection{\texttt{PUT /settings/search-preferences}}
Met à jour les préférences de recherche.

\subsubsection{\texttt{GET /settings/system}}
Retourne certains paramètres système visibles.

\subsection{Endpoints liés aux exécutions du système}

\subsubsection{\texttt{GET /runs}}
Liste les exécutions du système.

\subsubsection{\texttt{GET /runs/\{run\_id\}}}
Retourne le détail d’une exécution.

\subsubsection{\texttt{POST /runs/monitoring/trigger}}
Déclenche manuellement un cycle de veille.

\subsubsection{\texttt{POST /runs/mailbox-sync/trigger}}
Déclenche manuellement une synchronisation de messagerie.

\subsection{Règles générales de réponse}

Les règles suivantes s’appliquent aux réponses API :

\begin{itemize}
    \item toute réponse de lecture doit contenir les identifiants et états utiles ;
    \item toute réponse d’action doit préciser l’ancien état et le nouvel état si applicable ;
    \item les erreurs doivent être structurées ;
    \item les identifiants de journalisation doivent être retournés lorsqu’une action importante a été tracée.
\end{itemize}

\subsection{Périmètre minimal requis pour la version 1}

Le périmètre minimal d’API à rendre opérationnel dans la version 1 est le suivant :

\begin{itemize}
    \item consultation des offres ;
    \item détail d’une offre ;
    \item préparation d’un brouillon ;
    \item consultation d’un brouillon ;
    \item approbation d’un brouillon ;
    \item soumission d’un brouillon ;
    \item consultation des candidatures ;
    \item détail d’une candidature ;
    \item évaluation d’éligibilité à une relance ;
    \item préparation d’une relance.
\end{itemize}
\section{Spécification du frontend v1 : écrans, composants, données et actions utilisateur}

\subsection{Objectif}

Cette section décrit la structure fonctionnelle du frontend de la version 1.

Le frontend doit permettre à l’utilisateur de :

\begin{itemize}
    \item consulter les offres détectées ;
    \item comprendre les recommandations du système ;
    \item valider ou rejeter les brouillons préparés ;
    \item suivre les candidatures ;
    \item piloter les relances ;
    \item configurer les préférences du système.
\end{itemize}

\subsection{Principes généraux}

Le frontend de la version 1 suit les principes suivants :

\begin{itemize}
    \item interface orientée tableau de bord métier ;
    \item organisation centrée sur les objets fonctionnels ;
    \item lisibilité des états et des décisions ;
    \item priorité donnée à la validation et au suivi ;
    \item sobriété de l’interface.
\end{itemize}

\subsection{Écrans principaux}

Les écrans principaux du frontend v1 sont les suivants :

\begin{enumerate}
    \item tableau de bord ;
    \item inbox des offres ;
    \item détail d’une offre ;
    \item brouillons à valider ;
    \item pipeline des candidatures ;
    \item paramètres.
\end{enumerate}

\subsection{Tableau de bord}

Le tableau de bord fournit une vue synthétique du système.

Il doit afficher notamment :

\begin{itemize}
    \item le nombre de nouvelles offres ;
    \item le nombre d’offres qualifiées ;
    \item le nombre de brouillons à valider ;
    \item le nombre de candidatures en suivi ;
    \item le nombre de relances dues ;
    \item les dernières réponses reçues ;
    \item le dernier run du système.
\end{itemize}

\subsection{Inbox des offres}

L’inbox des offres permet de consulter les offres détectées et de les filtrer.

Chaque offre affichée doit présenter :

\begin{itemize}
    \item le titre ;
    \item l’entreprise ;
    \item la localisation ;
    \item la date ;
    \item le score ;
    \item l’état ;
    \item la recommandation ;
    \item la source.
\end{itemize}

Les principales actions sont :

\begin{itemize}
    \item ouvrir une offre ;
    \item rejeter une offre ;
    \item lancer une préparation de brouillon ;
    \item recalculer le score.
\end{itemize}

\subsection{Détail d’une offre}

L’écran de détail d’une offre doit afficher :

\begin{itemize}
    \item les informations générales de l’offre ;
    \item le score global et le score d’action ;
    \item la justification du matching ;
    \item les sous-scores par dimension ;
    \item la description normalisée ;
    \item l’historique résumé de l’offre.
\end{itemize}

Les actions principales sont :

\begin{itemize}
    \item préparer un brouillon ;
    \item recalculer le score ;
    \item rejeter l’offre.
\end{itemize}

\subsection{Brouillons à valider}

Cet écran est dédié au human-in-the-loop.

Il doit permettre de consulter pour chaque brouillon :

\begin{itemize}
    \item l’offre associée ;
    \item le CV sélectionné ;
    \item la lettre générée ;
    \item le message préparé ;
    \item les réponses préparées ;
    \item l’état du brouillon.
\end{itemize}

Les actions principales sont :

\begin{itemize}
    \item approuver ;
    \item rejeter ;
    \item demander des modifications ;
    \item reporter ;
    \item soumettre.
\end{itemize}

\subsection{Pipeline des candidatures}

Le pipeline des candidatures permet de suivre les candidatures envoyées.

Il doit afficher notamment :

\begin{itemize}
    \item le poste ;
    \item l’entreprise ;
    \item la date d’envoi ;
    \item le canal ;
    \item l’état courant ;
    \item la dernière activité ;
    \item les réponses reçues ;
    \item les relances éventuelles.
\end{itemize}

Les actions principales sont :

\begin{itemize}
    \item consulter une candidature ;
    \item consulter son historique ;
    \item évaluer l’éligibilité à une relance ;
    \item préparer une relance ;
    \item mettre à jour un statut.
\end{itemize}

\subsection{Paramètres}

L’écran des paramètres doit permettre de gérer :

\begin{itemize}
    \item les préférences de recherche ;
    \item le profil utilisateur ;
    \item les versions de CV ;
    \item la politique de relance.
\end{itemize}

\subsection{Composants principaux}

Les composants principaux à prévoir sont notamment :

\begin{itemize}
    \item des cartes KPI ;
    \item des badges d’état ;
    \item des badges de score ;
    \item des listes et tables d’offres ;
    \item des fiches de brouillons ;
    \item des vues de lettres et messages préparés ;
    \item des cartes de candidatures ;
    \item des formulaires de paramètres.
\end{itemize}

\subsection{Actions utilisateur minimales}

Les actions minimales à rendre disponibles dans la version 1 sont :

\begin{itemize}
    \item consulter une offre ;
    \item rejeter une offre ;
    \item préparer un brouillon ;
    \item approuver un brouillon ;
    \item rejeter un brouillon ;
    \item demander des modifications ;
    \item soumettre une candidature ;
    \item consulter une candidature ;
    \item préparer une relance ;
    \item modifier les préférences ;
    \item gérer les versions de CV.
\end{itemize}

\subsection{Navigation}

La navigation principale du frontend est organisée autour des rubriques suivantes :

\begin{itemize}
    \item tableau de bord ;
    \item offres ;
    \item brouillons ;
    \item candidatures ;
    \item relances ;
    \item paramètres.
\end{itemize}

\subsection{États UX à prévoir}

Chaque écran doit prévoir des états explicites de :

\begin{itemize}
    \item chargement ;
    \item absence de données ;
    \item erreur ;
    \item succès après action ;
    \item confirmation d’action sensible.
\end{itemize}

\subsection{Orientation générale}

Le frontend de la version 1 doit être conçu comme un poste de pilotage personnel, privilégiant la lisibilité, la rapidité d’exécution et le contrôle humain plutôt que la sophistication visuelle.

\section{Spécification des workflows frontend-backend v1}

\subsection{Objectif}

Cette section formalise les principaux workflows entre le frontend et le backend dans la version 1 du système.

L’objectif est de décrire, pour chaque parcours critique :

\begin{itemize}
    \item le point de départ dans l’interface ;
    \item l’action utilisateur ;
    \item l’appel API correspondant ;
    \item le traitement backend ;
    \item les transitions d’état éventuelles ;
    \item le retour visible dans l’interface ;
    \item la gestion des erreurs.
\end{itemize}

\subsection{Consultation des offres}

Depuis l’écran d’inbox des offres, l’utilisateur peut consulter la liste des offres détectées.

Ce workflow repose sur l’appel :

\begin{center}
\texttt{GET /offers}
\end{center}

Le backend retourne une liste paginée d’offres résumées.

Ce workflow n’implique aucun changement d’état métier.

\subsection{Consultation du détail d’une offre}

Depuis l’inbox, l’utilisateur peut ouvrir une offre particulière.

Ce workflow repose sur l’appel :

\begin{center}
\texttt{GET /offers/\{offer\_id\}}
\end{center}

Le backend retourne le détail de l’offre, son scoring, sa justification et son historique utile.

\subsection{Préparation d’un brouillon de candidature}

Depuis le détail d’une offre ou depuis l’inbox, l’utilisateur peut demander la préparation d’une candidature.

Ce workflow repose sur l’appel :

\begin{center}
\texttt{POST /offers/\{offer\_id\}/prepare-draft}
\end{center}

Le backend :

\begin{itemize}
    \item vérifie l’état de l’offre ;
    \item charge le profil et les préférences ;
    \item sélectionne un CV ;
    \item génère les éléments du brouillon ;
    \item crée un objet brouillon ;
    \item journalise la décision.
\end{itemize}

Ce workflow peut entraîner une transition d’état de l’offre et la création d’un brouillon prêt à validation.

\subsection{Consultation d’un brouillon}

Depuis l’écran des brouillons, l’utilisateur peut ouvrir un brouillon préparé.

Ce workflow repose sur l’appel :

\begin{center}
\texttt{GET /drafts/\{draft\_id\}}
\end{center}

Le backend retourne le détail complet du brouillon, y compris les artefacts préparés.

\subsection{Validation d’un brouillon}

Depuis le détail du brouillon, l’utilisateur peut :

\begin{itemize}
    \item approuver ;
    \item rejeter ;
    \item demander des modifications ;
    \item reporter la décision.
\end{itemize}

Ces workflows reposent sur les appels :

\begin{center}
\texttt{POST /drafts/\{draft\_id\}/approve}
\end{center}

\begin{center}
\texttt{POST /drafts/\{draft\_id\}/reject}
\end{center}

\begin{center}
\texttt{POST /drafts/\{draft\_id\}/request-changes}
\end{center}

\begin{center}
\texttt{POST /drafts/\{draft\_id\}/postpone}
\end{center}

Le backend applique la machine d’état, journalise la décision et retourne le nouvel état.

\subsection{Soumission d’une candidature}

Depuis un brouillon approuvé, l’utilisateur peut lancer la soumission.

Ce workflow repose sur l’appel :

\begin{center}
\texttt{POST /drafts/\{draft\_id\}/submit}
\end{center}

Le backend :

\begin{itemize}
    \item vérifie l’état du brouillon ;
    \item crée la candidature ;
    \item met à jour les états associés ;
    \item initialise le suivi ;
    \item journalise l’opération.
\end{itemize}

\subsection{Consultation du pipeline des candidatures}

Depuis l’écran pipeline, l’utilisateur peut consulter les candidatures suivies.

Ce workflow repose sur l’appel :

\begin{center}
\texttt{GET /applications}
\end{center}

Le backend retourne une liste filtrable et paginée des candidatures.

\subsection{Consultation du détail d’une candidature}

Depuis le pipeline, l’utilisateur peut ouvrir une candidature particulière.

Ce workflow repose sur l’appel :

\begin{center}
\texttt{GET /applications/\{application\_id\}}
\end{center}

Le backend retourne :

\begin{itemize}
    \item les informations générales ;
    \item le brouillon d’origine ;
    \item les réponses reçues ;
    \item les relances ;
    \item l’historique des décisions et transitions.
\end{itemize}

\subsection{Évaluation de l’éligibilité à une relance}

Depuis le détail d’une candidature, l’utilisateur peut demander si une relance est pertinente.

Ce workflow repose sur l’appel :

\begin{center}
\texttt{GET /applications/\{application\_id\}/followup-eligibility}
\end{center}

Le backend évalue la politique de relance et retourne une décision d’éligibilité.

\subsection{Préparation d’une relance}

Si la relance est éligible, l’utilisateur peut demander sa préparation.

Ce workflow repose sur l’appel :

\begin{center}
\texttt{POST /applications/\{application\_id\}/followups/prepare}
\end{center}

Le backend génère un brouillon de relance, le persiste et met à jour les états concernés.

\subsection{Validation et envoi d’une relance}

Une relance préparée peut ensuite être approuvée puis envoyée.

Ces workflows reposent sur les appels :

\begin{center}
\texttt{POST /followups/\{followup\_id\}/approve}
\end{center}

\begin{center}
\texttt{POST /followups/\{followup\_id\}/send}
\end{center}

Le backend applique les transitions d’état, journalise les décisions et met à jour la candidature associée.

\subsection{Modification des paramètres}

Depuis l’écran des paramètres, l’utilisateur peut modifier :

\begin{itemize}
    \item les préférences de recherche ;
    \item le profil utilisateur ;
    \item les versions de CV.
\end{itemize}

Ces workflows reposent notamment sur les appels :

\begin{center}
\texttt{PUT /settings/search-preferences}
\end{center}

\begin{center}
\texttt{PUT /profiles/active}
\end{center}

\begin{center}
\texttt{POST /profiles/active/cv-versions}
\end{center}

\subsection{Règles générales}

Les règles suivantes s’appliquent à tous les workflows frontend-backend :

\begin{itemize}
    \item toute action sensible doit retourner un résultat explicite ;
    \item toute transition d’état doit être visible dans l’interface ;
    \item les actions disponibles dans l’interface doivent être cohérentes avec l’état métier courant ;
    \item les erreurs doivent être affichées de manière compréhensible ;
    \item les listes doivent être réactualisées après les mutations importantes.
\end{itemize}

\subsection{Orientation pour la version 1}

La version 1 doit couvrir en priorité les parcours suivants :

\begin{itemize}
    \item consultation d’une offre ;
    \item préparation d’un brouillon ;
    \item validation humaine ;
    \item soumission d’une candidature ;
    \item consultation du pipeline ;
    \item préparation d’une relance.
\end{itemize}
\section{Stratégie de sécurité, confidentialité et gestion des accès v1}

\subsection{Objectif}

Cette section définit les principes de sécurité applicables à la version 1 du système.

L’objectif est de garantir :

\begin{itemize}
    \item la protection des données personnelles ;
    \item la protection des secrets techniques ;
    \item le contrôle des actions sensibles ;
    \item la traçabilité des opérations importantes ;
    \item une exposition minimale du système.
\end{itemize}

\subsection{Principes généraux}

Les principes suivants s’appliquent à la version 1 :

\begin{itemize}
    \item principe du moindre privilège ;
    \item validation humaine obligatoire avant toute action externe sensible ;
    \item séparation stricte entre secrets et code source ;
    \item minimisation des données stockées ;
    \item journalisation des actions critiques.
\end{itemize}

\subsection{Catégories de données sensibles}

Les catégories de données sensibles manipulées par le système comprennent notamment :

\begin{itemize}
    \item les données personnelles de l’utilisateur ;
    \item les CV et lettres de motivation ;
    \item les candidatures et leur historique ;
    \item les messages issus de la messagerie ;
    \item les secrets techniques et jetons d’accès ;
    \item les artefacts générés par le système.
\end{itemize}

\subsection{Gestion des secrets}

Les secrets techniques ne doivent jamais être stockés dans le code source.

Ils doivent être gérés via :

\begin{itemize}
    \item des variables d’environnement ;
    \item des fichiers de configuration locaux non versionnés ;
    \item un fichier \texttt{.env.example} sans données sensibles.
\end{itemize}

Aucun secret réel ne doit apparaître dans :

\begin{itemize}
    \item le dépôt Git ;
    \item les fichiers exportés ;
    \item les journaux applicatifs ;
    \item les captures de configuration.
\end{itemize}

\subsection{Gestion des accès à l’application}

La version 1 ne doit pas être exposée publiquement par défaut.

Les principes suivants sont retenus :

\begin{itemize}
    \item absence d’accès anonyme ;
    \item authentification minimale requise ;
    \item exposition réseau limitée ;
    \item protection des endpoints sensibles.
\end{itemize}

\subsection{Contrôle des actions sensibles}

Les actions suivantes sont considérées comme sensibles :

\begin{itemize}
    \item soumettre une candidature ;
    \item envoyer une relance ;
    \item modifier des paramètres critiques ;
    \item connecter ou synchroniser la messagerie ;
    \item déclencher certaines automatisations externes.
\end{itemize}

Ces actions doivent :

\begin{itemize}
    \item passer par le backend ;
    \item être validées vis-à-vis de l’état métier courant ;
    \item être journalisées ;
    \item rester soumises au contrôle humain lorsque nécessaire.
\end{itemize}

\subsection{Messagerie et intégrations externes}

Les intégrations externes doivent respecter le principe du moindre privilège.

En particulier :

\begin{itemize}
    \item les permissions d’accès à la messagerie doivent être limitées ;
    \item les synchronisations doivent être tracées ;
    \item les envois doivent rester contrôlés ;
    \item les données récupérées doivent être limitées à ce qui est utile.
\end{itemize}

\subsection{Protection des fichiers}

Les fichiers générés ou utilisés par le système doivent être stockés dans un espace non public.

Les règles suivantes s’appliquent :

\begin{itemize}
    \item stockage privé ;
    \item noms de fichiers contrôlés ;
    \item pas d’exposition directe non sécurisée ;
    \item conservation des versions importantes ;
    \item possibilité d’archivage ou de suppression.
\end{itemize}

\subsection{Politique de logs et d’audit}

Le système doit journaliser :

\begin{itemize}
    \item les exécutions importantes ;
    \item les transitions d’état ;
    \item les validations utilisateur ;
    \item les soumissions et relances ;
    \item les erreurs techniques significatives.
\end{itemize}

En revanche, les journaux ne doivent pas contenir en clair :

\begin{itemize}
    \item mots de passe ;
    \item jetons ;
    \item secrets ;
    \item détails sensibles non nécessaires.
\end{itemize}

\subsection{Gestion des erreurs}

Les erreurs exposées à l’utilisateur doivent rester compréhensibles sans divulguer d’informations techniques sensibles.

Les détails techniques complets doivent être conservés dans les journaux internes, et non dans les réponses frontend.

\subsection{Intégrité métier}

Les règles métier suivantes participent à la sécurité logique du système :

\begin{itemize}
    \item un brouillon non approuvé ne peut pas être soumis ;
    \item une candidature rejetée, acceptée ou clôturée ne peut pas être relancée ;
    \item une transition d’état interdite ne peut pas être contournée ;
    \item les états terminaux ne peuvent pas être quittés librement.
\end{itemize}

\subsection{Sauvegarde et résilience minimale}

La version 1 doit prévoir un minimum de sauvegarde des éléments suivants :

\begin{itemize}
    \item base de données ;
    \item fichiers utilisateur ;
    \item préférences ;
    \item historique de candidatures ;
    \item workflows exportés.
\end{itemize}

\subsection{Orientation générale}

La version 1 ne vise pas une architecture de sécurité avancée de type multi-tenant ou entreprise complète.

Elle doit néanmoins reposer sur des règles simples mais strictes, permettant d’obtenir un système personnel sérieux, non exposé inutilement et respectueux des données manipulées.
\section{Stratégie d’observabilité, logs, monitoring et debug v1}

\subsection{Objectif}

Cette section définit les principes d’observabilité de la version 1 du système.

L’objectif est de garantir que le système soit :

\begin{itemize}
    \item compréhensible ;
    \item diagnosable ;
    \item mesurable ;
    \item traçable ;
    \item débogable.
\end{itemize}

\subsection{Principes généraux}

Les principes suivants s’appliquent à la version 1 :

\begin{itemize}
    \item tout workflow important doit être traçable ;
    \item les logs doivent être structurés ;
    \item les logs techniques et les logs métier doivent être distingués ;
    \item toute exécution importante doit disposer d’un identifiant unique ;
    \item les erreurs ne doivent pas être silencieuses.
\end{itemize}

\subsection{Zones à observer}

Les zones principales à observer sont les suivantes :

\begin{itemize}
    \item les exécutions du système ;
    \item les connecteurs de collecte ;
    \item les décisions métier ;
    \item les transitions d’état ;
    \item les intégrations externes ;
    \item les performances générales.
\end{itemize}

\subsection{Politique de logs}

Le système doit utiliser au minimum les niveaux de logs suivants :

\begin{itemize}
    \item \texttt{DEBUG} ;
    \item \texttt{INFO} ;
    \item \texttt{WARNING} ;
    \item \texttt{ERROR}.
\end{itemize}

Les logs doivent inclure autant que possible :

\begin{itemize}
    \item un horodatage ;
    \item un niveau ;
    \item un composant ;
    \item un identifiant de run ;
    \item un type d’événement ;
    \item un identifiant d’objet concerné ;
    \item un message explicite ;
    \item des métadonnées utiles.
\end{itemize}

\subsection{Traces d’exécution}

Chaque exécution importante du système doit être historisée.

Cela concerne notamment :

\begin{itemize}
    \item les runs de veille ;
    \item les synchronisations de messagerie ;
    \item les runs de scoring ;
    \item les préparations de candidatures ;
    \item les traitements de relance.
\end{itemize}

Chaque run doit être associé à :

\begin{itemize}
    \item un identifiant unique ;
    \item un type de run ;
    \item une date de début ;
    \item une date de fin ;
    \item un statut ;
    \item des compteurs utiles ;
    \item un résumé.
\end{itemize}

\subsection{Événements métier à journaliser}

Les événements métier suivants doivent être journalisés :

\begin{itemize}
    \item détection d’une offre ;
    \item normalisation d’une offre ;
    \item décision de déduplication ;
    \item scoring d’une offre ;
    \item création d’un brouillon ;
    \item approbation ou rejet d’un brouillon ;
    \item soumission d’une candidature ;
    \item détection d’une réponse ;
    \item préparation ou envoi d’une relance ;
    \item début et fin d’un run système.
\end{itemize}

\subsection{Métriques utiles}

La version 1 doit permettre le suivi de métriques simples, notamment :

\begin{itemize}
    \item nombre d’offres collectées ;
    \item nombre d’offres nouvelles ;
    \item nombre d’offres qualifiées ;
    \item nombre de brouillons générés ;
    \item nombre de candidatures soumises ;
    \item nombre de réponses reçues ;
    \item nombre de relances préparées ;
    \item nombre d’erreurs critiques ;
    \item temps moyen d’exécution des runs.
\end{itemize}

\subsection{Monitoring minimal}

La version 1 doit prévoir un monitoring minimal comprenant :

\begin{itemize}
    \item des endpoints de santé ;
    \item l’accès aux journaux applicatifs ;
    \item l’historique des runs ;
    \item une vue synthétique de l’état du système ;
    \item une visibilité sur les erreurs récentes.
\end{itemize}

\subsection{Stratégie de debug}

Le debug doit être possible à trois niveaux :

\begin{itemize}
    \item par exécution du système ;
    \item par objet métier ;
    \item par intégration externe.
\end{itemize}

Le système doit permettre de répondre à des questions telles que :

\begin{itemize}
    \item quel run a échoué ;
    \item sur quel objet ;
    \item à quelle étape ;
    \item pour quelle raison ;
    \item avec quel impact métier.
\end{itemize}

\subsection{Aides au diagnostic}

La version 1 doit prévoir au minimum :

\begin{itemize}
    \item des identifiants de corrélation ;
    \item des journaux structurés ;
    \item l’historique des décisions ;
    \item des déclenchements manuels contrôlés pour certains workflows ;
    \item une consultation des historiques d’objets et de runs.
\end{itemize}

\subsection{Orientation générale}

L’observabilité de la version 1 doit rester simple mais rigoureuse.

Elle doit privilégier :

\begin{itemize}
    \item la compréhension du comportement du système ;
    \item la capacité de diagnostic ;
    \item la traçabilité des décisions ;
    \item la correction rapide des erreurs.
\end{itemize}
\section{Stratégie de test, validation et qualité logicielle v1}

\subsection{Objectif}

Cette section définit la stratégie de test de la version 1 du système.

L’objectif est de garantir :

\begin{itemize}
    \item la fiabilité des modules ;
    \item la cohérence des workflows métier ;
    \item la stabilité du système lors de son évolution ;
    \item la détection rapide des erreurs.
\end{itemize}

\subsection{Niveaux de tests}

La stratégie de test repose sur plusieurs niveaux :

\begin{itemize}
    \item tests unitaires ;
    \item tests d’intégration ;
    \item tests de workflow métier ;
    \item tests de régression ;
    \item tests manuels contrôlés.
\end{itemize}

\subsection{Tests unitaires}

Les tests unitaires vérifient le fonctionnement isolé des modules.

Les modules prioritaires à tester sont :

\begin{itemize}
    \item la normalisation des offres ;
    \item la déduplication ;
    \item le calcul des scores ;
    \item la sélection de CV ;
    \item la politique de relance.
\end{itemize}

\subsection{Tests d’intégration}

Les tests d’intégration vérifient la coopération entre modules.

Ils couvrent notamment :

\begin{itemize}
    \item l’ingestion complète d’une offre ;
    \item la préparation d’un brouillon de candidature ;
    \item la soumission d’une candidature.
\end{itemize}

\subsection{Tests de workflow métier}

Les tests de workflow vérifient les parcours utilisateur complets.

Les workflows principaux à tester sont :

\begin{itemize}
    \item offre détectée vers candidature soumise ;
    \item candidature soumise vers relance.
\end{itemize}

\subsection{Tests de régression}

Les tests de régression permettent de vérifier que les modifications du code n’introduisent pas de régression.

Ils doivent être exécutés avant toute nouvelle version importante.

\subsection{Tests manuels}

Certains éléments nécessitent des tests manuels :

\begin{itemize}
    \item qualité des lettres générées ;
    \item comportement des connecteurs web ;
    \item intégration avec la messagerie ;
    \item cohérence globale de l’interface.
\end{itemize}

\subsection{Jeux de données de test}

La version 1 doit inclure des jeux de données de test comprenant :

\begin{itemize}
    \item des offres fictives ;
    \item des candidatures simulées ;
    \item des réponses email simulées.
\end{itemize}

\subsection{Tests des machines d’état}

Les transitions d’état doivent être testées pour :

\begin{itemize}
    \item les offres ;
    \item les brouillons ;
    \item les candidatures ;
    \item les relances.
\end{itemize}

Les transitions interdites doivent être explicitement vérifiées.

\subsection{Validation avant déploiement}

Avant toute nouvelle version, les vérifications suivantes doivent être effectuées :

\begin{itemize}
    \item exécution des tests unitaires ;
    \item exécution des tests d’intégration ;
    \item validation du workflow principal ;
    \item vérification des logs ;
    \item vérification de la cohérence de l’interface.
\end{itemize}

\subsection{Orientation générale}

La version 1 privilégie une stratégie de test simple mais rigoureuse, centrée sur les modules critiques et les workflows métier principaux.

\section{Roadmap de développement et planification des versions}

\subsection{Objectif}

Cette section définit la planification du développement du système.

L’objectif est de structurer la construction du produit en plusieurs versions progressives.

\subsection{Vision générale}

Le développement est organisé en quatre étapes :

\begin{itemize}
\item MVP (produit minimal fonctionnel)
\item version 1 (copilote complet)
\item version 2 (automatisation intelligente)
\item version 3 (agents autonomes avancés)
\end{itemize}

\subsection{MVP}

Le MVP vise à produire un système minimal mais immédiatement utile.

Les fonctionnalités principales sont :

\begin{itemize}
\item collecte d’offres depuis quelques sources
\item normalisation simple
\item scoring basique
\item inbox d’offres
\item préparation de brouillons de candidature
\item validation humaine
\item suivi simple des candidatures
\end{itemize}

\subsection{Version 1}

La version 1 constitue la première version complète du système.

Elle ajoute notamment :

\begin{itemize}
\item déduplication robuste
\item matching avancé
\item justification des scores
\item pipeline complet de candidatures
\item politique de relance
\item synchronisation avec la messagerie
\item observabilité et journalisation complètes
\end{itemize}

\subsection{Version 2}

La version 2 introduit des capacités d’automatisation intelligente.

Les fonctionnalités supplémentaires comprennent :

\begin{itemize}
\item apprentissage des préférences utilisateur
\item sélection automatique du CV
\item enrichissement des informations d’entreprise
\item priorisation intelligente des offres
\item alertes proactives
\end{itemize}

\subsection{Version 3}

La version 3 introduit des agents autonomes spécialisés.

Ces agents peuvent :

\begin{itemize}
\item explorer les sources d’offres
\item analyser automatiquement les nouvelles opportunités
\item préparer des candidatures optimisées
\item ajuster la stratégie de recherche
\end{itemize}

\subsection{Ordre de développement recommandé}

L’ordre de développement recommandé est :

\begin{enumerate}
\item architecture backend et base de données
\item ingestion et normalisation des offres
\item matching et scoring
\item interface frontend minimale
\item préparation des brouillons
\item validation humaine
\item pipeline de candidatures
\item gestion des relances
\item intégration messagerie
\item observabilité et monitoring
\end{enumerate}

\subsection{Orientation générale}

Le développement doit privilégier une progression incrémentale, chaque version apportant une valeur fonctionnelle réelle tout en préparant les capacités plus avancées des versions ultérieures.

\section{Plan de démarrage concret du projet}

\subsection{Objectif}

Cette section a pour objectif de transformer le cahier des charges en plan de démarrage concret du projet.

Elle précise :

\begin{itemize}
    \item la structure initiale du dépôt ;
    \item les premiers composants à mettre en place ;
    \item l’ordre de construction recommandé ;
    \item les premiers incréments fonctionnels à viser.
\end{itemize}

\subsection{Structure initiale du projet}

Le projet est initialisé sous la forme d’un mono-dépôt contenant au minimum :

\begin{itemize}
    \item un backend ;
    \item un frontend ;
    \item un worker ;
    \item des workflows ;
    \item une documentation ;
    \item des données de test ;
    \item des fichiers d’infrastructure.
\end{itemize}

\subsection{Fichiers racine initiaux}

Les fichiers suivants doivent être créés dès le démarrage :

\begin{itemize}
    \item \texttt{README.md} ;
    \item \texttt{.gitignore} ;
    \item \texttt{.env.example} ;
    \item \texttt{docker-compose.yml} ;
    \item \texttt{Makefile}.
\end{itemize}

\subsection{Premier objectif technique}

Le premier objectif technique recommandé est le suivant :

\begin{center}
être capable d’enregistrer une offre en base, de l’exposer via l’API et de l’afficher dans une interface simple.
\end{center}

Cet objectif permet de valider simultanément :

\begin{itemize}
    \item la base de données ;
    \item le backend ;
    \item l’API ;
    \item le frontend ;
    \item la circulation des données.
\end{itemize}

\subsection{Premier incrément fonctionnel}

Le premier incrément fonctionnel recommandé consiste à mettre en place :

\begin{itemize}
    \item une table \texttt{offers} ;
    \item un modèle d’offre ;
    \item un repository d’offres ;
    \item des endpoints de lecture ;
    \item une page frontend de liste des offres ;
    \item une page frontend de détail d’une offre ;
    \item un jeu de données initial.
\end{itemize}

\subsection{Incréments fonctionnels suivants}

Les incréments suivants sont recommandés dans l’ordre suivant :

\begin{enumerate}
    \item scoring des offres ;
    \item préparation d’un brouillon ;
    \item validation et soumission ;
    \item pipeline des candidatures ;
    \item préparation des relances.
\end{enumerate}

\subsection{Ordre de développement recommandé}

L’ordre de développement concret recommandé est :

\begin{itemize}
    \item mise en place du noyau technique du backend ;
    \item mise en place du modèle d’offre ;
    \item mise en place de l’API des offres ;
    \item mise en place du frontend minimal ;
    \item ajout du scoring ;
    \item ajout des brouillons ;
    \item ajout des candidatures ;
    \item ajout des relances.
\end{itemize}

\subsection{Premier sprint recommandé}

Le premier sprint vise à produire une fondation fonctionnelle.

Son objectif est :

\begin{center}
afficher une liste d’offres fictives depuis PostgreSQL dans l’interface.
\end{center}

Les livrables du sprint 1 sont :

\begin{itemize}
    \item backend exécutable ;
    \item base PostgreSQL connectée ;
    \item migration initiale ;
    \item seed d’offres ;
    \item endpoints de lecture ;
    \item interface de liste et détail des offres.
\end{itemize}

\subsection{Orientation générale}

Le démarrage du projet doit privilégier des incréments verticaux simples, validant rapidement la cohérence entre base de données, backend, API et frontend, avant d’introduire les couches d’automatisation et d’agentique plus avancées.

\section{Liste exacte des premiers fichiers du Sprint 1}

\subsection{Objectif}

Cette section précise les premiers fichiers à créer pour démarrer concrètement le projet.

L’objectif du Sprint 1 est de rendre possible l’affichage d’offres fictives stockées en base dans une interface web simple.

\subsection{Fichiers racine}

Les fichiers suivants doivent être créés à la racine du projet :

\begin{itemize}
    \item \texttt{README.md} ;
    \item \texttt{.gitignore} ;
    \item \texttt{.env.example} ;
    \item \texttt{docker-compose.yml} ;
    \item \texttt{Makefile}.
\end{itemize}

\subsection{Fichiers backend}

Les fichiers backend à créer en priorité sont :

\begin{itemize}
    \item \texttt{backend/pyproject.toml} ;
    \item \texttt{backend/alembic.ini} ;
    \item \texttt{backend/app/main.py} ;
    \item \texttt{backend/app/core/config.py} ;
    \item \texttt{backend/app/core/database.py} ;
    \item \texttt{backend/app/core/logging.py} ;
    \item \texttt{backend/app/domain/enums/offer\_state.py} ;
    \item \texttt{backend/app/infrastructure/orm/models/offer.py} ;
    \item \texttt{backend/app/schemas/offer.py} ;
    \item \texttt{backend/app/repositories/offer\_repository.py} ;
    \item \texttt{backend/app/services/offer\_service.py} ;
    \item \texttt{backend/app/api/router.py} ;
    \item \texttt{backend/app/api/routes/health.py} ;
    \item \texttt{backend/app/api/routes/offers.py} ;
    \item \texttt{backend/app/api/deps.py}.
\end{itemize}

\subsection{Fichiers de migration et de données}

Les fichiers suivants doivent être prévus :

\begin{itemize}
    \item une migration initiale dans \texttt{backend/migrations/} ;
    \item un fichier de seed d’offres ;
    \item un script de chargement des données initiales.
\end{itemize}

\subsection{Fichiers frontend}

Les fichiers frontend à créer en priorité sont :

\begin{itemize}
    \item \texttt{frontend/package.json} ;
    \item \texttt{frontend/tsconfig.json} ;
    \item \texttt{frontend/next.config.ts} ;
    \item \texttt{frontend/src/app/layout.tsx} ;
    \item \texttt{frontend/src/app/page.tsx} ;
    \item \texttt{frontend/src/app/offers/page.tsx} ;
    \item \texttt{frontend/src/app/offers/[id]/page.tsx} ;
    \item \texttt{frontend/src/components/layout/AppShell.tsx} ;
    \item \texttt{frontend/src/components/offers/OfferList.tsx} ;
    \item \texttt{frontend/src/components/offers/OfferCard.tsx} ;
    \item \texttt{frontend/src/components/offers/OfferDetail.tsx} ;
    \item \texttt{frontend/src/components/ui/StatusBadge.tsx} ;
    \item \texttt{frontend/src/components/ui/ScoreBadge.tsx} ;
    \item \texttt{frontend/src/lib/api-client.ts} ;
    \item \texttt{frontend/src/features/offers/api.ts} ;
    \item \texttt{frontend/src/features/offers/types.ts}.
\end{itemize}

\subsection{Fichiers de test}

Les fichiers de test minimaux recommandés sont :

\begin{itemize}
    \item \texttt{backend/tests/test\_health.py} ;
    \item \texttt{backend/tests/test\_offer\_repository.py} ;
    \item \texttt{backend/tests/test\_offers\_api.py}.
\end{itemize}

\subsection{Ordre de création recommandé}

L’ordre de création recommandé est :

\begin{enumerate}
    \item fichiers racine ;
    \item noyau technique du backend ;
    \item modèle d’offre ;
    \item repository et service d’offre ;
    \item endpoints API ;
    \item migration initiale et seed ;
    \item frontend minimal ;
    \item tests minimaux.
\end{enumerate}

\subsection{Définition de fini du Sprint 1}

Le Sprint 1 est considéré comme terminé lorsque :

\begin{itemize}
    \item le backend est exécutable ;
    \item la base PostgreSQL est connectée ;
    \item la table des offres existe ;
    \item des offres fictives sont chargées en base ;
    \item l’API permet de lister les offres ;
    \item l’API permet de consulter le détail d’une offre ;
    \item le frontend affiche la liste des offres ;
    \item le frontend affiche le détail d’une offre.
\end{itemize}

\section{Backlog détaillé du Sprint 1}

\subsection{Objectif}

Le Sprint 1 a pour objectif de rendre possible l’affichage d’offres fictives stockées dans PostgreSQL dans une interface frontend simple.

\subsection{Découpage du Sprint 1}

Le Sprint 1 est structuré en six blocs :

\begin{enumerate}
    \item initialisation du dépôt ;
    \item socle backend ;
    \item modèle Offre et persistance ;
    \item API des offres ;
    \item frontend minimal ;
    \item tests et validation.
\end{enumerate}

\subsection{Bloc 1 : initialisation du dépôt}

Les tâches de ce bloc sont :

\begin{itemize}
    \item créer le dépôt et l’arborescence initiale ;
    \item créer les fichiers racine ;
    \item préparer l’environnement local.
\end{itemize}

\subsection{Bloc 2 : socle backend}

Les tâches de ce bloc sont :

\begin{itemize}
    \item initialiser le projet Python backend ;
    \item créer le point d’entrée FastAPI ;
    \item mettre en place la configuration ;
    \item mettre en place la connexion à la base ;
    \item initialiser le système de migrations.
\end{itemize}

\subsection{Bloc 3 : modèle Offre et persistance}

Les tâches de ce bloc sont :

\begin{itemize}
    \item définir l’énumération minimale des états d’offre ;
    \item créer le modèle ORM de l’offre ;
    \item créer la migration initiale ;
    \item créer les schémas d’offre ;
    \item créer le repository d’offre ;
    \item créer le service d’offre ;
    \item préparer les données de seed ;
    \item créer le script de chargement initial.
\end{itemize}

\subsection{Bloc 4 : API des offres}

Les tâches de ce bloc sont :

\begin{itemize}
    \item créer les dépendances API ;
    \item créer la route de santé ;
    \item créer les routes d’offre ;
    \item agréger les routes dans le routeur principal.
\end{itemize}

\subsection{Bloc 5 : frontend minimal}

Les tâches de ce bloc sont :

\begin{itemize}
    \item initialiser le projet Next.js ;
    \item créer le layout global ;
    \item créer le client API frontend ;
    \item définir les types frontend d’offre ;
    \item définir les appels API liés aux offres ;
    \item créer les composants UI minimaux ;
    \item créer la page liste des offres ;
    \item créer la page détail d’une offre ;
    \item créer la page d’accueil minimale.
\end{itemize}

\subsection{Bloc 6 : tests et validation}

Les tâches de ce bloc sont :

\begin{itemize}
    \item tester l’endpoint de santé ;
    \item tester le repository d’offre ;
    \item tester l’API des offres ;
    \item effectuer une validation manuelle de bout en bout.
\end{itemize}

\subsection{Ordre recommandé}

L’ordre recommandé d’exécution est :

\begin{enumerate}
    \item initialisation du dépôt ;
    \item socle backend ;
    \item modèle Offre et persistance ;
    \item API des offres ;
    \item frontend minimal ;
    \item tests et validation.
\end{enumerate}

\subsection{Tâches critiques}

Les tâches critiques du Sprint 1 sont :

\begin{itemize}
    \item connexion à la base de données ;
    \item création du modèle d’offre ;
    \item migration initiale ;
    \item création du repository d’offre ;
    \item création des endpoints d’offre ;
    \item chargement des données fictives ;
    \item affichage de la liste des offres ;
    \item affichage du détail d’une offre.
\end{itemize}

\subsection{Définition de fini du Sprint 1}

Le Sprint 1 est considéré comme terminé lorsque :

\begin{itemize}
    \item la stack est exécutable ;
    \item PostgreSQL est connecté ;
    \item la table des offres existe ;
    \item des offres fictives sont chargées en base ;
    \item l’API permet de lister les offres ;
    \item l’API permet de consulter le détail d’une offre ;
    \item le frontend affiche la liste des offres ;
    \item le frontend affiche le détail d’une offre.
\end{itemize}

\section{Décisions techniques structurantes}

Cette section fixe les décisions techniques fondamentales qui structurent le développement du système.

Ces décisions servent de référence pendant toute la phase d’implémentation et visent à garantir la cohérence de l’architecture ainsi que la stabilité du développement.

\subsection{Architecture initiale du système}

La première version du système adopte une architecture simple de type monolithe modulaire.

Le système est composé des éléments suivants :

\begin{itemize}
\item une interface utilisateur web (frontend) ;
\item une API backend responsable de la logique métier ;
\item une base de données PostgreSQL ;
\item un ou plusieurs workers dédiés aux tâches asynchrones.
\end{itemize}

L’utilisation de microservices indépendants n’est pas retenue dans les premières versions afin de limiter la complexité opérationnelle et de favoriser un développement rapide et cohérent.

\subsection{Source de vérité métier}

La logique métier et les états des objets principaux du système sont centralisés dans le backend.

Les états des entités suivantes doivent être contrôlés exclusivement par le backend :

\begin{itemize}
\item les offres ;
\item les brouillons de candidature ;
\item les candidatures ;
\item les relances.
\end{itemize}

Le frontend, les agents ou les outils d’orchestration ne doivent pas modifier directement ces états sans passer par les services du backend.

\subsection{Stratégie de collecte des offres}

La collecte des offres est organisée par niveaux de complexité.

\begin{itemize}
\item \textbf{Niveau 1 :} pages carrières d’entreprises et sources simples ;
\item \textbf{Niveau 2 :} plateformes publiques d’offres d’emploi ;
\item \textbf{Niveau 3 :} plateformes complexes nécessitant authentification ou navigation avancée.
\end{itemize}

Les premières versions du système privilégient les sources de niveau 1 afin de valider le fonctionnement global du pipeline avant d’intégrer des sources plus complexes.

\subsection{Stratégie de scoring des offres}

L’évaluation de la pertinence d’une offre repose sur un score global combinant plusieurs critères.

Ces critères peuvent inclure :

\begin{itemize}
\item la pertinence métier ;
\item la correspondance des compétences ;
\item la localisation ;
\item le type de contrat ;
\item le délai de candidature ;
\item l’intérêt stratégique de l’entreprise.
\end{itemize}

Le score doit rester explicable et permettre de justifier les décisions prises par le système.

\subsection{Niveau d’automatisation autorisé}

Le système vise à assister l’utilisateur tout en conservant un contrôle humain sur les actions externes.

Les principes suivants s’appliquent :

\begin{itemize}
\item la préparation des candidatures peut être automatisée ;
\item la soumission d’une candidature nécessite une validation humaine ;
\item les relances doivent être proposées par le système mais validées par l’utilisateur ;
\item aucune candidature ne doit être envoyée automatiquement sans approbation explicite.
\end{itemize}

Ce principe garantit la qualité des candidatures et la confiance de l’utilisateur dans le système.


\end{document}





